% Options for packages loaded elsewhere
\PassOptionsToPackage{unicode}{hyperref}
\PassOptionsToPackage{hyphens}{url}
\PassOptionsToPackage{dvipsnames,svgnames,x11names}{xcolor}
\documentclass[
  11pt,
]{article}
\usepackage{xcolor}
\usepackage[margin=1in]{geometry}
\usepackage{amsmath,amssymb}
\setcounter{secnumdepth}{-\maxdimen} % remove section numbering
\usepackage{iftex}
\ifPDFTeX
  \usepackage[T1]{fontenc}
  \usepackage[utf8]{inputenc}
  \usepackage{textcomp} % provide euro and other symbols
\else % if luatex or xetex
  \usepackage{unicode-math} % this also loads fontspec
  \defaultfontfeatures{Scale=MatchLowercase}
  \defaultfontfeatures[\rmfamily]{Ligatures=TeX,Scale=1}
\fi
\usepackage{lmodern}
\ifPDFTeX\else
  % xetex/luatex font selection
\fi
% Use upquote if available, for straight quotes in verbatim environments
\IfFileExists{upquote.sty}{\usepackage{upquote}}{}
\IfFileExists{microtype.sty}{% use microtype if available
  \usepackage[]{microtype}
  \UseMicrotypeSet[protrusion]{basicmath} % disable protrusion for tt fonts
}{}
\makeatletter
\@ifundefined{KOMAClassName}{% if non-KOMA class
  \IfFileExists{parskip.sty}{%
    \usepackage{parskip}
  }{% else
    \setlength{\parindent}{0pt}
    \setlength{\parskip}{6pt plus 2pt minus 1pt}}
}{% if KOMA class
  \KOMAoptions{parskip=half}}
\makeatother
\usepackage{color}
\usepackage{fancyvrb}
\newcommand{\VerbBar}{|}
\newcommand{\VERB}{\Verb[commandchars=\\\{\}]}
\DefineVerbatimEnvironment{Highlighting}{Verbatim}{commandchars=\\\{\}}
% Add ',fontsize=\small' for more characters per line
\newenvironment{Shaded}{}{}
\newcommand{\AlertTok}[1]{\textcolor[rgb]{1.00,0.00,0.00}{\textbf{#1}}}
\newcommand{\AnnotationTok}[1]{\textcolor[rgb]{0.38,0.63,0.69}{\textbf{\textit{#1}}}}
\newcommand{\AttributeTok}[1]{\textcolor[rgb]{0.49,0.56,0.16}{#1}}
\newcommand{\BaseNTok}[1]{\textcolor[rgb]{0.25,0.63,0.44}{#1}}
\newcommand{\BuiltInTok}[1]{\textcolor[rgb]{0.00,0.50,0.00}{#1}}
\newcommand{\CharTok}[1]{\textcolor[rgb]{0.25,0.44,0.63}{#1}}
\newcommand{\CommentTok}[1]{\textcolor[rgb]{0.38,0.63,0.69}{\textit{#1}}}
\newcommand{\CommentVarTok}[1]{\textcolor[rgb]{0.38,0.63,0.69}{\textbf{\textit{#1}}}}
\newcommand{\ConstantTok}[1]{\textcolor[rgb]{0.53,0.00,0.00}{#1}}
\newcommand{\ControlFlowTok}[1]{\textcolor[rgb]{0.00,0.44,0.13}{\textbf{#1}}}
\newcommand{\DataTypeTok}[1]{\textcolor[rgb]{0.56,0.13,0.00}{#1}}
\newcommand{\DecValTok}[1]{\textcolor[rgb]{0.25,0.63,0.44}{#1}}
\newcommand{\DocumentationTok}[1]{\textcolor[rgb]{0.73,0.13,0.13}{\textit{#1}}}
\newcommand{\ErrorTok}[1]{\textcolor[rgb]{1.00,0.00,0.00}{\textbf{#1}}}
\newcommand{\ExtensionTok}[1]{#1}
\newcommand{\FloatTok}[1]{\textcolor[rgb]{0.25,0.63,0.44}{#1}}
\newcommand{\FunctionTok}[1]{\textcolor[rgb]{0.02,0.16,0.49}{#1}}
\newcommand{\ImportTok}[1]{\textcolor[rgb]{0.00,0.50,0.00}{\textbf{#1}}}
\newcommand{\InformationTok}[1]{\textcolor[rgb]{0.38,0.63,0.69}{\textbf{\textit{#1}}}}
\newcommand{\KeywordTok}[1]{\textcolor[rgb]{0.00,0.44,0.13}{\textbf{#1}}}
\newcommand{\NormalTok}[1]{#1}
\newcommand{\OperatorTok}[1]{\textcolor[rgb]{0.40,0.40,0.40}{#1}}
\newcommand{\OtherTok}[1]{\textcolor[rgb]{0.00,0.44,0.13}{#1}}
\newcommand{\PreprocessorTok}[1]{\textcolor[rgb]{0.74,0.48,0.00}{#1}}
\newcommand{\RegionMarkerTok}[1]{#1}
\newcommand{\SpecialCharTok}[1]{\textcolor[rgb]{0.25,0.44,0.63}{#1}}
\newcommand{\SpecialStringTok}[1]{\textcolor[rgb]{0.73,0.40,0.53}{#1}}
\newcommand{\StringTok}[1]{\textcolor[rgb]{0.25,0.44,0.63}{#1}}
\newcommand{\VariableTok}[1]{\textcolor[rgb]{0.10,0.09,0.49}{#1}}
\newcommand{\VerbatimStringTok}[1]{\textcolor[rgb]{0.25,0.44,0.63}{#1}}
\newcommand{\WarningTok}[1]{\textcolor[rgb]{0.38,0.63,0.69}{\textbf{\textit{#1}}}}
\usepackage{longtable,booktabs,array}
\usepackage{calc} % for calculating minipage widths
% Correct order of tables after \paragraph or \subparagraph
\usepackage{etoolbox}
\makeatletter
\patchcmd\longtable{\par}{\if@noskipsec\mbox{}\fi\par}{}{}
\makeatother
% Allow footnotes in longtable head/foot
\IfFileExists{footnotehyper.sty}{\usepackage{footnotehyper}}{\usepackage{footnote}}
\makesavenoteenv{longtable}
\setlength{\emergencystretch}{3em} % prevent overfull lines
\providecommand{\tightlist}{%
  \setlength{\itemsep}{0pt}\setlength{\parskip}{0pt}}
% DDC v1.0.3 — pandoc header for pdfLaTeX. AuthorityEffect: None.
\usepackage{newunicodechar}
\newunicodechar{Λ}{\ensuremath{\Lambda}}
\newunicodechar{≠}{\ensuremath{\neq}}
\newunicodechar{−}{\ensuremath{-}}
\newunicodechar{Φ}{\ensuremath{\Phi}}
\newunicodechar{Γ}{\ensuremath{\Gamma}}
\newunicodechar{ℓ}{\ensuremath{\ell}}
\newunicodechar{₁}{\ensuremath{_1}}
\newunicodechar{₂}{\ensuremath{_2}}
\usepackage{pdflscape}
\emergencystretch=3em
\usepackage{microtype}
% Let long identifiers in code spans break after underscores.
\let\ddcunderscore\_
\renewcommand{\_}{\ddcunderscore\allowbreak}
\ifx\pdftrailerid\undefined\else\pdftrailerid{ddc-1.0.3}\fi
\AtBeginDocument{\hypersetup{
  pdftitle={Discrete Deformation Calculus: A Finite Trace Calculus for Obstruction Handling},
  pdfauthor={Jeremy H. Carroll},
  pdfsubject={Version 1.0.3; 2026-10-03; preprint; DOI pending},
  pdfkeywords={discrete deformation calculus, finite obstruction calculus, trace calculus, Lean 4}
}}
\usepackage{bookmark}
\IfFileExists{xurl.sty}{\usepackage{xurl}}{} % add URL line breaks if available
\urlstyle{same}
\hypersetup{
  colorlinks=true,
  linkcolor={Maroon},
  filecolor={Maroon},
  citecolor={Blue},
  urlcolor={Blue},
  pdfcreator={LaTeX via pandoc}}

\author{}
\date{}

\begin{document}

\begin{center}
\phantomsection\label{discrete-deformation-calculus}
{\LARGE\bfseries Discrete Deformation Calculus}\par\medskip
{\large A Finite Trace Calculus for Obstruction Handling}\par\medskip
Jeremy H. Carroll\par
Version 1.0.3 --- 2026-10-03\par
Text: CC BY 4.0 --- Code: PolyForm Noncommercial 1.0.0
\end{center}

\section{Abstract}\label{abstract}

The Discrete Deformation Calculus (DDC) records what a finite
computation did to a declared obstruction. Each step emits one letter
from \(\mathcal{A}_{\mathrm{DDC}} = \{E, R, C, I, \Lambda, D, M, T\}\):
exact repair, residual projection, context exposure, integration,
licensed lift, descent, typed refusal and terminalization. A step
carries a receipt, and a tape is the hash-chained sequence of receipts.
Diagnostics come from CODA-FOC: \(\Phi_1\) is the weighted \(\ell^1\)
distance of a residual to the image of \(d_0\), and
\(\Gamma_2 = \lVert d_1 x \rVert_1\) is its face-closure defect.

This version states the letter laws as a finite-control transition
system over unbounded data and proves its tape results in Lean 4 without
Mathlib. Every proved or computed claim is bound to a Lean declaration,
a Python test or a written proof shipped in this package; the rest are
typed as open problems or positioning. The main results are these. Only
a descent with a loss ledger or a licensed lift can lower \(\Phi_1\) or
\(\Gamma_2\). Consequently an open obstruction (\(\Gamma_2 > 0\)) never
reaches a terminal without such a priced move. The tape C·R·E·D·I·T is
lawful exactly when each of its six steps satisfies its law. A lawful
one starts from \(\Phi_1 = \Gamma_2 = 0\), and without a loss ledger at
D it is not lawful. Every step the paper admits satisfies the deposited
historical comparison laws on the same abstract state; no
deployed-runtime equivalence is claimed. A lawful tape is replay
evidence; it is never authority.

\section{0. Scope and nonclaims}\label{scope-and-nonclaims}

\begin{Shaded}
\begin{Highlighting}[]
\NormalTok{DDC does not prove that reality is discrete or that the universe is a cochain complex.}
\NormalTok{DDC does not prove quantum mechanics, continuum physics, consciousness, semantic truth,}
\NormalTok{autonomous self{-}modification, or global topology.}

\NormalTok{DDC is not authority. A DDC tape is replay evidence; a receipt is bounded verification}
\NormalTok{data; a terminal is a routing terminal. No authorization is conferred.}

\NormalTok{"All problems are obstruction problems" is a programme thesis, not a theorem.}
\NormalTok{"The universe is not very creative" is motivation, not runtime mathematics.}
\end{Highlighting}
\end{Shaded}

Every diagnostic in this paper is taken on a declared finite carrier
with a declared detector. Claims about trace equivalence, normal forms,
detector adequacy in general, and lift minimality remain open (§13).

\section{Part I --- The calculus}\label{part-i-the-calculus}

\section{1. Definition}\label{definition}

\textbf{Definition 1.1} (OE010-09). A \emph{DDC regime} is a finite
carrier, declared coboundary maps \(d_0, d_1\) with \(d_1 d_0 = 0\), a
detector that measures \(\Phi_1\) and \(\Gamma_2\), and a terminal
policy. A \emph{step} is a letter from \(\mathcal{A}_{\mathrm{DDC}}\),
the state after it, and the witnesses it carries. A \emph{tape} is a
finite sequence of steps; it is \emph{lawful} when every step satisfies
its letter's step law (§7) and the control laws (§10).

Carrier, ZeroRel, Observer, Witness, Residual, Quotient, CompletionLaw
and AuthorityEffect are used in their CODA sense; this paper relies only
on the ones it defines in this section and in §3 (OE010-02).

The letters classify the role a step played in the life of an
obstruction:

\begin{Shaded}
\begin{Highlighting}[]
\NormalTok{context exposure · exact repair · residual projection · integration}
\NormalTok{licensed lift · descent · typed refusal · terminalization}
\end{Highlighting}
\end{Shaded}

\section{2. Core invariant}\label{core-invariant}

DDC does not only classify obstruction states. It classifies what each
computation did to an obstruction:

\[\boxed{\text{DDC is a typed calculus of what a computation did to an obstruction.}}\]

A status system says \emph{pass}, \emph{fail} or \emph{unknown}. DDC
asks which finite operation occurred: did it expose a carrier, repair a
defect, project a residual, integrate a receipt, license a lift, descend
with loss, refuse, or terminalize? This lets error handling be compared
across domains without claiming that the domains share an ontology.

\section{3. Cochain window and
diagnostics}\label{cochain-window-and-diagnostics}

When a cochain regime is declared, every execution window is

\[C^0 \xrightarrow{d_0} C^1 \xrightarrow{d_1} C^2, \qquad d_1 \circ d_0 = 0.\]

\begin{longtable}[]{@{}ll@{}}
\toprule\noalign{}
Object & Runtime reading \\
\midrule\noalign{}
\endhead
\bottomrule\noalign{}
\endlastfoot
\(C^0\) & vertex variables, potentials, repair parameters \\
\(C^1\) & edge defects, claim gaps, local residuals \\
\(C^2\) & face and loop closure defects \\
\end{longtable}

The convention is adapter-conditional: a prompt, a patch or an agent
trace does not have \(C^\bullet\), \(d_0\), \(d_1\), \(\Phi_1\) or
\(\Gamma_2\) until a carrier, a comparison map, a repair image and a
detector have been declared.

\textbf{Definition 3.1} (FOC diagnostics). For a residual \(x \in C^1\)
with positive edge weights \(w\),
\(\Phi_1(x) = \min_{\varphi} \sum_e w_e \lvert x_e - (d_0\varphi)_e \rvert\)
and \(\Gamma_2(x) = \lVert d_1 x \rVert_1\). The class of \(x\) is
\emph{ActiveClosureFailure} when \(\Gamma_2 > 0\), otherwise
\emph{ExactRepair} when \(\Phi_1 = 0\), otherwise \emph{ClosedResidual}.
A positive \(\Gamma_2\) dominates (witness dominance).

\textbf{Proposition 3.2} (DDC-TAPE-11). On a finite graph, \(\Phi_1\) is
attained at a potential determined by a spanning forest whose edges are
all tight. Enumerating spanning forests therefore computes \(\Phi_1\)
exactly.

\emph{Proof.} Fix one vertex potential to zero in each connected
component. Positive weights bound each edge difference on an objective
sublevel set; paths in a spanning forest then bound every remaining
potential. Apply compactness to the real extension; a minimum is
attained. The objective is piecewise linear on finitely many polyhedral
regions. The minimizing face in one such region is nonempty and bounded,
so it contains a vertex of that region. At such a vertex, independent
tight constraints \(x_e = \varphi_v - \varphi_u\) fix \(\varphi\).
Independent incidence rows form a forest, and fixing \(\varphi\) needs
one spanning each component. A minimizing vertex determined by a
spanning forest has rational coordinates, so the optimum is also
attained over the rationals. \(\square\)

\texttt{cochain.\allowbreak phi1} implements the enumeration and refuses as
\texttt{BUDGET\_UNKNOWN} above a declared subset budget;
\texttt{test\_cochain.\allowbreak py::test\_phi1\_is\_gauge\_invariant\_and\_matches\_brute\_force\_on\_small\_grids}
checks it against an integer search. \texttt{cochain.\allowbreak phi1\_certificate}
also returns a dual witness: a flow \(y\) with \(d_0^{\top} y = 0\) and
\(\lvert y_e \rvert \le w_e\) whose value \(\langle y, x \rangle\)
equals the primal objective. Dual vertices set \(y_e = \pm w_e\) off a
spanning forest and solve the forest edges by conservation, so the same
enumeration finds them, and strong duality is checked rather than
assumed
(\texttt{test\_cochain.\allowbreak py::test\_phi1\_certificates\_close\_the\_duality\_gap\_exactly}).
Every \(\Phi_1\) in the demos is measured through this certificate.

A fixed forest representative need not minimize this objective. For the
complete graph on \(a<b<c<d\), orient all six edges increasingly and put
\(-1\) on each. The forest rooted at \(a\) leaves the three other edges
at \(-1\), so the recorded residual norm is 3, while \(\Phi_1=2\). Both
representatives have the same quotient class. Chapter 2 distinguishes
these quantities.

\textbf{Theorem 3.3} (OE010-01, DDC-TAPE-10). Over \(\mathbb{Q}\) with
the standard inner product,
\(C^1 = \operatorname{im} d_0 \oplus \operatorname{im} d_1^{\top}
\oplus \mathcal{H}^1\) with
\(\mathcal{H}^1 = \ker d_1 \cap \ker d_0^{\top}\). Write
\(x = x_E + x_C + x_H\). Then \(\Gamma_2(x) > 0\) exactly when
\(x_C \neq 0\), and \(\Phi_1(x) = 0\) exactly when \(x_C = 0\) and
\(x_H = 0\).

\emph{Proof.} Set \(U=\operatorname{im}d_0\) and
\(V=\operatorname{im}d_1^{\top}\). The chain law gives \(U\perp V\), and
\((U\oplus V)^\perp=\ker d_0^{\top}\cap\ker d_1=\mathcal H^1\).
Finite-dimensional positive definiteness gives the unique orthogonal
direct sum \(C^1=U\oplus V\oplus\mathcal H^1\). The defining linear
systems have rational coefficients, so their unique projected vectors
are rational. Now \(d_1\) kills \(x_E\) because \(d_1 d_0 = 0\) and
kills \(x_H\) by definition. On \(\operatorname{im} d_1^{\top}\) it is
injective: \(d_1 d_1^{\top}\psi = 0\) implies
\(\lVert d_1^{\top}\psi \rVert^2 = \langle \psi, d_1 d_1^{\top}\psi
\rangle = 0\). So \(d_1 x = d_1 x_C\) vanishes only when \(x_C = 0\).
\(\Phi_1(x) = 0\) means \(x \in \operatorname{im} d_0\), which by the
direct sum means \(x_C = x_H = 0\). \(\square\)

The subscripts name Hodge components, not letters. The letter E acts on
the exact component; the co-exact component is what \(\Gamma_2\) sees;
the harmonic component is the closed residual.
\texttt{cochain.\allowbreak hodge\_split} computes the split exactly and
\texttt{test\_cochain.\allowbreak py::test\_hodge\_split\_is\_orthogonal\_and\_matches\_the\_classes}
checks the theorem on four complexes.

\section{4. The alphabet}\label{the-alphabet}

Every state transition or structural modification in a DDC regime emits
exactly one letter.

\begin{longtable}[]{@{}
  >{\raggedright\arraybackslash}p{(\linewidth - 6\tabcolsep) * \real{0.1304}}
  >{\raggedright\arraybackslash}p{(\linewidth - 6\tabcolsep) * \real{0.1739}}
  >{\raggedright\arraybackslash}p{(\linewidth - 6\tabcolsep) * \real{0.3478}}
  >{\raggedright\arraybackslash}p{(\linewidth - 6\tabcolsep) * \real{0.3478}}@{}}
\toprule\noalign{}
\begin{minipage}[b]{\linewidth}\raggedright
Letter
\end{minipage} & \begin{minipage}[b]{\linewidth}\raggedright
Name
\end{minipage} & \begin{minipage}[b]{\linewidth}\raggedright
Runtime meaning
\end{minipage} & \begin{minipage}[b]{\linewidth}\raggedright
Must not claim
\end{minipage} \\
\midrule\noalign{}
\endhead
\bottomrule\noalign{}
\endlastfoot
\texttt{E} & Exact repair & removes a defect that lies in
\(\operatorname{im} d_0\) & semantic truth, global correctness,
authority \\
\texttt{R} & Residual projection & records the residual class; keeps
\(\Phi_1\) and \(\Gamma_2\) & resolution \\
\texttt{C} & Context exposure & declares or splits the carrier, window
or boundary & repair or authority \\
\texttt{I} & Integration & appends a receipt to the hash-chained tape &
promotion \\
\texttt{Λ} & Structural lift & extends the carrier under a licence &
arbitrary growth or automatic promotion \\
\texttt{D} & Descent & re-indexes or coarsens with a declared loss
ledger & lossless transport unless witnessed \\
\texttt{M} & Typed refusal & refuses an invalid, unwitnessed or
out-of-budget request & failure of the product \\
\texttt{T} & Terminalization & closes the tape on a routing terminal &
authority \\
\end{longtable}

\begin{Shaded}
\begin{Highlighting}[]
\NormalTok{DDC records how an obstruction was handled, not whether the world is true.}
\end{Highlighting}
\end{Shaded}

\section{5. Step receipts}\label{step-receipts}

Every admitted abstract step produces a deterministic record of the
supplied states, carrier labels and witness flags, linked to the
previous record. These flags are hypotheses about evidence presence;
they do not authenticate a receipt, a loss ledger, a licence or detector
adequacy. Direct construction of \texttt{State} likewise does not
certify its diagnostics.

\texttt{measure} computes exact diagnostics. Passing a carrier map to
\texttt{run\_tape} recomputes the initial and successor diagnostics,
checks that R changes the cochain by a coboundary, and binds topology
and positive weights into the tape genesis digest. Without that map,
replay checks only the abstract laws. The Lean theorems concern those
laws, not an authenticated runtime. The record is:

\begin{Shaded}
\begin{Highlighting}[]
\FunctionTok{\{}
  \DataTypeTok{"step\_index"}\FunctionTok{:} \DecValTok{2}\FunctionTok{,}
  \DataTypeTok{"before\_digest"}\FunctionTok{:} \StringTok{"sha256 of the canonical before{-}state"}\FunctionTok{,}
  \DataTypeTok{"after\_digest"}\FunctionTok{:} \StringTok{"sha256 of the canonical after{-}state"}\FunctionTok{,}
  \DataTypeTok{"carrier\_before"}\FunctionTok{:} \StringTok{"two\_triangles"}\FunctionTok{,}
  \DataTypeTok{"carrier\_after"}\FunctionTok{:} \StringTok{"two\_triangles"}\FunctionTok{,}
  \DataTypeTok{"class\_before"}\FunctionTok{:} \StringTok{"ExactRepair"}\FunctionTok{,}
  \DataTypeTok{"class\_after"}\FunctionTok{:} \StringTok{"ExactRepair"}\FunctionTok{,}
  \DataTypeTok{"phi\_before"}\FunctionTok{:} \StringTok{"0"}\FunctionTok{,} \DataTypeTok{"phi\_after"}\FunctionTok{:} \StringTok{"0"}\FunctionTok{,}
  \DataTypeTok{"gamma\_before"}\FunctionTok{:} \StringTok{"0"}\FunctionTok{,} \DataTypeTok{"gamma\_after"}\FunctionTok{:} \StringTok{"0"}\FunctionTok{,}
  \DataTypeTok{"witness"}\FunctionTok{:} \FunctionTok{\{}
    \DataTypeTok{"receipt"}\FunctionTok{:} \KeywordTok{false}\FunctionTok{,}
    \DataTypeTok{"loss\_ledger"}\FunctionTok{:} \KeywordTok{false}\FunctionTok{,}
    \DataTypeTok{"lift\_licence"}\FunctionTok{:} \KeywordTok{false}\FunctionTok{,}
    \DataTypeTok{"reason"}\FunctionTok{:} \KeywordTok{null}
  \FunctionTok{\},}
  \DataTypeTok{"control\_after"}\FunctionTok{:} \StringTok{"Open"}\FunctionTok{,}
  \DataTypeTok{"routing\_terminal"}\FunctionTok{:} \KeywordTok{null}\FunctionTok{,}
  \DataTypeTok{"authority\_effect"}\FunctionTok{:} \StringTok{"None"}
\FunctionTok{\}}
\end{Highlighting}
\end{Shaded}

The integration tape chains receipts as
\(h_{t+1} = \mathrm{SHA256}(t : h_t : \alpha : \mathrm{json}(r_t))\)
with canonical JSON (sorted keys, exact rationals as strings, floats
refused). Production receipts add a trace identifier, detector and
terminalizer versions, protected distinctions, and references to loss
ledgers and lift requests. No receipt stores raw user text by default.
Tamper evidence assumes the collision resistance of SHA-256 and an
independently retained head and length (DDC-I-02). Unanchored replay
checks internal consistency only; a rewritten or truncated history can
remain internally consistent.

\section{6. Diagnostic and runtime
crosswalk}\label{diagnostic-and-runtime-crosswalk}

\begin{longtable}[]{@{}
  >{\raggedright\arraybackslash}p{(\linewidth - 6\tabcolsep) * \real{0.2353}}
  >{\raggedright\arraybackslash}p{(\linewidth - 6\tabcolsep) * \real{0.1765}}
  >{\raggedright\arraybackslash}p{(\linewidth - 6\tabcolsep) * \real{0.2353}}
  >{\raggedright\arraybackslash}p{(\linewidth - 6\tabcolsep) * \real{0.3529}}@{}}
\toprule\noalign{}
\begin{minipage}[b]{\linewidth}\raggedright
DDC term
\end{minipage} & \begin{minipage}[b]{\linewidth}\raggedright
Runtime term
\end{minipage} & \begin{minipage}[b]{\linewidth}\raggedright
Diagnostic meaning
\end{minipage} & \begin{minipage}[b]{\linewidth}\raggedright
Routing
\end{minipage} \\
\midrule\noalign{}
\endhead
\bottomrule\noalign{}
\endlastfoot
closure defect & \(\Gamma_2\) & \(d_1 x \neq 0\) & \texttt{Repair} or
\texttt{Refusal}; never \texttt{Check\allowbreak Passed} \\
irreducible residual & \(\Phi_1\) & quotient mass, \(\Phi_1 > 0\) &
\texttt{Requires\allowbreak Lift} or residual terminal \\
exact class & \texttt{Exact\allowbreak Repair} & \(\Phi_1 = 0\), \(\Gamma_2 = 0\) &
\texttt{Check\allowbreak Passed} within the declared scope \\
open defect & \texttt{Active\allowbreak Closure\allowbreak Failure} & \(\Gamma_2 > 0\) &
\texttt{Repair} or \texttt{Refusal} \\
stabilized survivor & \texttt{Closed\allowbreak Residual} & \(\Phi_1 > 0\),
\(\Gamma_2 = 0\) & residual terminal or \texttt{Requires\allowbreak Lift} \\
\end{longtable}

The four axes stay separate:

\begin{Shaded}
\begin{Highlighting}[]
\NormalTok{DiagnosticClass ≠ RoutingTerminal ≠ ExtensionStatus ≠ AuthorityStatus.}
\end{Highlighting}
\end{Shaded}

Letters are not terminals. \texttt{Check\allowbreak Passed} on the routing axis is
not a licence: this package adopts \texttt{F\_term:\ REFUSED} as its
explicit unlicensed-terminal policy. This is a policy definition, not an
application of an external theorem.

\section{7. The laws}\label{the-laws}

The eight candidate laws of the first draft are now step and control
laws of a finite-control transition system (Definition 10.1). The
comparison columns below refer exclusively to the deposited
\texttt{engine\allowbreak Law} and \texttt{engine\_violations} specifications,
retained from repository commit
\texttt{e4a1214bb3e47570b0f74ecc9117c98f27a26b5e}. Their names are
historical. No theorem here establishes parity with a deployed Rust
implementation; claims about later runtime changes require a separate
pinned conformance check.

\begin{enumerate}
\def\labelenumi{\arabic{enumi}.}
\tightlist
\item
  \textbf{Trace typing} (OE010-03). Every letter carries its input
  class, output class, required witness and refusal condition, and every
  lawful step satisfies its letter's step law
  (\texttt{Ddc.\allowbreak well\_typed\_step}).
\item
  \textbf{Terminal absorption} (OE010-10). After \texttt{T} no letter is
  admitted, so \texttt{T} is the last letter of a lawful tape
  (\texttt{Ddc.\allowbreak lawful\_T\_is\_last}).
\item
  \textbf{Refusal absorption unless lifted} (DDC-TAPE-01). After
  \texttt{M} only \texttt{C} (changed scope), \texttt{Λ} (licensed lift)
  or \texttt{T} may follow.
\item
  \textbf{Lift trigger.} \texttt{Λ} needs a licence and a trigger: the
  state before it is a \texttt{Closed\allowbreak Residual} or its detector is
  inadequate. A trigger is not a licence.
\item
  \textbf{Integration requires a receipt.} \texttt{I} without a receipt
  is refused.
\item
  \textbf{Descent requires a loss ledger.} \texttt{D} without a loss
  ledger is refused.
\item
  \textbf{Exact repair discharges.} \texttt{E} needs an exact class
  before it, leaves the carrier unchanged, and leaves
  \(\Phi_1 = \Gamma_2 = 0\) and a zero defect after it. A gauge repair
  cannot cross classes (Chapter 1), so \texttt{E} is never a disguised
  class change.
\item
  \textbf{Context exposure is not repair.} \texttt{C} changes the
  carrier and never lowers \(\Phi_1\) or \(\Gamma_2\). Hiding an
  obstruction by re-scoping has to be priced as \texttt{D} or
  \texttt{Λ}.
\end{enumerate}

The step laws, compared with the deposited historical specification:

\begin{longtable}[]{@{}
  >{\raggedright\arraybackslash}p{(\linewidth - 6\tabcolsep) * \real{0.1000}}
  >{\raggedright\arraybackslash}p{(\linewidth - 6\tabcolsep) * \real{0.3000}}
  >{\raggedright\arraybackslash}p{(\linewidth - 6\tabcolsep) * \real{0.3000}}
  >{\raggedright\arraybackslash}p{(\linewidth - 6\tabcolsep) * \real{0.3000}}@{}}
\toprule\noalign{}
\begin{minipage}[b]{\linewidth}\raggedright
Letter
\end{minipage} & \begin{minipage}[b]{\linewidth}\raggedright
Step law
\end{minipage} & \begin{minipage}[b]{\linewidth}\raggedright
Deposited comparison
\end{minipage} & \begin{minipage}[b]{\linewidth}\raggedright
Additional paper law
\end{minipage} \\
\midrule\noalign{}
\endhead
\bottomrule\noalign{}
\endlastfoot
\texttt{E} & same carrier; after \(\Phi_1 = \Gamma_2 = 0\) & yes
(\texttt{Exact\allowbreak Repair\allowbreak Left\allowbreak Residual}) & before class exact; zero defect
after; adequate detector \\
\texttt{R} & same carrier; \(\Phi_1\) unchanged & yes
(\texttt{Residual\allowbreak Projection\allowbreak Changed\allowbreak Mass}) & \(\Gamma_2\) unchanged;
adequate detector \\
\texttt{C} & carrier changes & yes
(\texttt{Carrier\allowbreak Unchanged\allowbreak Under\allowbreak Restructure}) & \(\Phi_1\) and
\(\Gamma_2\) do not fall (Law 8) \\
\texttt{I} & state unchanged & yes (\texttt{Inert\allowbreak Op\allowbreak Mutated\allowbreak State}) &
receipt present (Law 5) \\
\texttt{Λ} & carrier changes & yes & licence and trigger (Law 4) \\
\texttt{D} & carrier changes & yes & loss ledger (Law 6) \\
\texttt{M} & state unchanged & yes & reason from the closed
vocabulary \\
\texttt{T} & state unchanged; \(\Gamma_2 = 0\) & yes
(\texttt{Terminalized\allowbreak Open\allowbreak Residual}) & adequate detector \\
\end{longtable}

\textbf{Theorem 7.1} (DDC-TAPE-03). Every step admitted by the paper
laws satisfies the deposited comparison laws, represented on the
abstract state (\texttt{engine\allowbreak Law} in Lean, \texttt{engine\_violations}
in Python). The comparison is defined by the local source bytes; the
theorem does not execute Rust or assert correspondence with its current
version. Lean: \texttt{Ddc.\allowbreak paper\_refines\_engine}; Python:
\texttt{test\_automaton.\allowbreak py::test\_paper\_laws\_refine\_the\_engine\_laws\_on\_every\_step}.

\section{8. Five minimal demos}\label{five-minimal-demos}

Each demo is measured on a declared complex and replayed through the
laws by \texttt{code/\allowbreak ddc/\allowbreak demos.\allowbreak py}. Unlawful variants are kept as
falsifiers.
\texttt{test\_fsm\_and\_demos.\allowbreak py::test\_every\_demo\_reaches\_its\_stated\_verdict}
pins every verdict below.

\textbf{Demo 1 --- boundary cancellation} (two triangles sharing an
edge). An exact defect is registered, repaired and sealed:
\texttt{C·R·E·T}, routing \texttt{Check\allowbreak Passed}. A defect with
circulation around one face (\(\Gamma_2 = 1\)) ends in a typed refusal:
\texttt{C·R·M}. Its falsifier \texttt{C·R·E·T} is refused at \texttt{E},
because the class before \texttt{E} is not exact.

\textbf{Demo 2 --- local green, global red} (a filled triangle with
value 1 on every edge). Each edge is locally consistent, yet
\(d_1 x = 3\), so \(\Gamma_2 = 3\) and \(\Phi_1 = 3\). The lawful tape
is \texttt{C·R·M}. Appending \texttt{T} is refused
(\texttt{Terminalized\allowbreak Open\allowbreak Residual}): a refusal does not license a seal.

\textbf{Demo 3 --- stabilized survivor} (a hollow square with a unit
defect on one edge; \(\Phi_1 = 1\), \(\Gamma_2 = 0\)). \texttt{C·R·T}
seals on the residual terminal. A licensed lift that adds capacity away
from the cycle, \texttt{C·R·Λ·T}, re-measures \(\Phi_1 = 1\): the
survivor persists, as Chapter 5 proves it must. Without a licence the
record is \texttt{C·R·M·T} with \texttt{REFUSED\_REQUIRES\_LIFT}. Two
falsifiers: a lift without a licence is refused, and a lift that
attaches a face bounding the cycle turns the survivor into
\(\Gamma_2 = 1\), after which \texttt{T} is refused.

\textbf{Demo 4 --- detector collapse} (a constant detector reports zero
on the survivor). The lawful tape is \texttt{C·M} with
\texttt{REFUSED\_DETECTOR\_INADEQUATE}; \texttt{C·M·T} is refused
(\texttt{Blind\allowbreak Detector}), because a collapsing detector is how false
green enters. A licensed detector lift, \texttt{C·Λ·T}, re-measures
\(\Phi_1 = 1\) and seals honestly on the residual terminal.

\textbf{Demo 5 --- descent with loss} (a path of four edges carrying
\((1, 2, \tfrac{3}{2}, \tfrac{1}{4})\), coarsened to two edges). With a
loss ledger the tape \texttt{D·I·T} books loss \(\tfrac{7}{4}\) and
seals. Without one, \texttt{D} is refused (\texttt{Missing\allowbreak Loss\allowbreak Ledger}),
and the lawful record is \texttt{M·T} with
\texttt{REFUSED\_MISSING\_LOSS\_LEDGER}.

\section{9. Runtime context and verification
ceiling}\label{runtime-context-and-verification-ceiling}

The Rust crate \texttt{omega\_discrete\_deformation\_calculus} is
related runtime context. Names in \texttt{docs/\allowbreak claim\_code\_map.\allowbreak yaml}
are informational pointers, not self-contained evidence of current
engine conformance. This deposit proves refinement only of its own
historical comparison specification. It supplies no authenticated
licence evaluator and no test of the deployed Rust system.

\section{Part II --- The eight letters}\label{part-ii-the-eight-letters}

\begin{longtable}[]{@{}
  >{\raggedright\arraybackslash}p{(\linewidth - 4\tabcolsep) * \real{0.1667}}
  >{\raggedright\arraybackslash}p{(\linewidth - 4\tabcolsep) * \real{0.1667}}
  >{\raggedright\arraybackslash}p{(\linewidth - 4\tabcolsep) * \real{0.6667}}@{}}
\toprule\noalign{}
\begin{minipage}[b]{\linewidth}\raggedright
Chapter
\end{minipage} & \begin{minipage}[b]{\linewidth}\raggedright
Letter
\end{minipage} & \begin{minipage}[b]{\linewidth}\raggedright
Central result
\end{minipage} \\
\midrule\noalign{}
\endhead
\bottomrule\noalign{}
\endlastfoot
\hyperref[chapter-1]{1} & \texttt{E} & gauge repair discharges an exact
defect and cannot cross classes \\
\hyperref[chapter-2]{2} & \texttt{R} & projection keeps the class, with
one coordinate per independent cycle \\
\hyperref[chapter-3]{3} & \texttt{C} & exposure changes the carrier and
never lowers the diagnostics \\
\hyperref[chapter-4]{4} & \texttt{I} & the receipt tape only grows and
every letter enters the chain \\
\hyperref[chapter-5]{5} & \texttt{Λ} & a lift preserves norms; an
extending lift cannot exactify a non-exact residual \\
\hyperref[chapter-6]{6} & \texttt{D} & a descent is lossless exactly
when the dropped tail is zero \\
\hyperref[chapter-7]{7} & \texttt{M} & a refusal leaves the state
unchanged and names a reason \\
\hyperref[chapter-8]{8} & \texttt{T} & the abstract seal needs
\(\Gamma_2 = 0\) and adequate detection; continuation policy is a
separate guard \\
\end{longtable}

\section{Part III --- Tapes}\label{part-iii-tapes}

\section{10. The tape automaton}\label{the-tape-automaton}

\textbf{Definition 10.1.} A tape runs from a control state
\texttt{Open}, \texttt{Refused} or \texttt{Terminalized} and an abstract
data state. A step is admitted when the control admits its letter and
the letter's step law holds. \texttt{T} moves the control to
\texttt{Terminalized}, \texttt{M} to \texttt{Refused}, and \texttt{C} or
\texttt{Λ} after a refusal back to \texttt{Open}. \texttt{Terminalized}
admits nothing; \texttt{Refused} admits only \texttt{C}, \texttt{Λ} and
\texttt{T}. Replay stops at the first unlawful step, which is not
recorded. The Python automaton is \texttt{code/\allowbreak ddc/\allowbreak automaton.\allowbreak py}; the
Lean one is \texttt{Ddc.\allowbreak run} in \texttt{lean4/\allowbreak Ddc/\allowbreak Law.\allowbreak lean}. They state
the same laws. The Lean model represents supplied \(\Phi_1, \Gamma_2\)
by \(\mathbb{N}\) values; the laws use only equality, order and zero
tests, which the non-negative rationals of the Python model share.

\section{11. Tape theorems}\label{tape-theorems}

Every theorem below is proved in \texttt{lean4/\allowbreak } with no axiom outside
\{\texttt{propext}, \texttt{Quot.\allowbreak sound}\}
(\texttt{lean4/\allowbreak axiom\_report.\allowbreak txt}) and checked on examples in
\texttt{code/\allowbreak tests/\allowbreak }.

\textbf{Theorem 11.1} (OE010-03, composition). If replay of \(t_1\)
succeeds, replaying \(t_1 t_2\) equals replaying \(t_2\) from that
successful endpoint. If \(t_1\) fails, the concatenated replay fails as
well. The empty tape is the identity (\texttt{Ddc.\allowbreak run\_append},
\texttt{Ddc.\allowbreak lawful\_append}, \texttt{Ddc.\allowbreak lawful\_nil}). Lawful tapes
therefore form the morphisms of a category whose objects are
control-state pairs.

\textbf{Theorem 11.2} (OE010-10, terminal absorption). No step is
admitted after \texttt{T} (\texttt{Ddc.\allowbreak terminal\_absorbing},
\texttt{Ddc.\allowbreak lawful\_T\_is\_last}).

\textbf{Theorem 11.3} (DDC-TAPE-01, refusal absorption). After
\texttt{M} only \texttt{C}, \texttt{Λ} or \texttt{T} is admitted
(\texttt{Ddc.\allowbreak refusal\_absorbing}).

\textbf{Theorem 11.4} (DDC-TAPE-02, well-typed steps). Every step of a
lawful tape satisfies its letter's step law at the state it was taken
from (\texttt{Ddc.\allowbreak well\_typed\_step}).

\textbf{Theorem 11.5} (OE010-04, no authority). Every receipt a replay
emits carries \texttt{Authority\allowbreak Effect.\allowbreak none}, although the type has a
second value (\texttt{Ddc.\allowbreak authority\_effect\_none}); the receipt tape
only grows (\texttt{Ddc.\allowbreak receipts\_prefix}).

\textbf{Theorem 11.6} (OE010-17, no free decrease). If a lawful step
lowers \(\Phi_1\) or \(\Gamma_2\), its letter is \texttt{D} or
\texttt{Λ} (\texttt{Ddc.\allowbreak no\_free\_decrease}).

\emph{Proof sketch.} \texttt{E} starts from \(\Phi_1 = \Gamma_2 = 0\);
\texttt{R} keeps both; \texttt{C} cannot lower either; \texttt{I},
\texttt{M} and \texttt{T} leave the state unchanged. \(\square\)

\textbf{Theorem 11.7} (DDC-TAPE-04, open obstructions need a priced
move). From a state with \(\Gamma_2 > 0\) and a non-terminal control, a
lawful tape with no \texttt{D} and no \texttt{Λ} keeps \(\Gamma_2 > 0\)
and never terminalizes (\texttt{Ddc.\allowbreak open\_needs\_priced\_move}). This is
witness dominance at the level of tapes: an open obstruction reaches a
seal only through a booked loss or a licensed lift.
\texttt{test\_automaton.\allowbreak py::test\_open\_obstruction\_is\_never\_sealed\_without\_a\_priced\_move}
checks it by exhaustive reachability on a 32-state abstraction.

\textbf{Theorem 11.8} (CREDIT). The mnemonic ``CREDIT'' names the tape
C·R·E·D·I·T; it is only a mnemonic.

\begin{itemize}
\tightlist
\item
  (DDC-TAPE-05) CREDIT is lawful iff each of its six steps satisfies its
  law (\texttt{Ddc.\allowbreak credit\_lawful\_iff}).
\item
  (DDC-TAPE-06) A lawful CREDIT starts from \(\Phi_1 = \Gamma_2 = 0\)
  (\texttt{Ddc.\allowbreak credit\_starts\_exact}): \texttt{E} needs an exact class,
  \texttt{R} keeps the diagnostics and \texttt{C} cannot lower them, so
  the residual was exact before \texttt{C}.
\item
  (DDC-TAPE-07) Without a loss ledger at \texttt{D} it is not lawful
  (\texttt{Ddc.\allowbreak credit\_needs\_ledger}).
\item
  (DDC-TAPE-08) It is realizable (\texttt{Ddc.\allowbreak credit\_realizable};
  Python: \texttt{demos.\allowbreak credit}, an exact defect on a hollow triangle
  repaired and then booked down to one edge).
\end{itemize}

CREDIT is therefore the exact-repair tape with a booked descent. It is
not a route for dissolving an obstruction.

\textbf{Theorem 11.9} (OE010-16, driver soundness). The §12 driver never
proposes \texttt{T} on an open or blind state, proposes \texttt{E} only
on an exact class, and proposes \texttt{Λ} only with a Law 4 trigger
(\texttt{Ddc.\allowbreak driver\_never\_seals\_open},
\texttt{Ddc.\allowbreak driver\_repairs\_only\_exact},
\texttt{Ddc.\allowbreak driver\_lift\_trigger}).

\section{12. The §23 driver}\label{the-23-driver}

The first draft generated tapes with a state machine whose table swapped
two classes: it tagged \texttt{Active\allowbreak Closure\allowbreak Failure} with \(\Phi_1 > 0\)
and \texttt{Closed\allowbreak Residual} with \(\Gamma_2 > 0\). The corrected driver
(OE010-16) proposes:

\begin{longtable}[]{@{}
  >{\raggedright\arraybackslash}p{(\linewidth - 2\tabcolsep) * \real{0.5000}}
  >{\raggedright\arraybackslash}p{(\linewidth - 2\tabcolsep) * \real{0.5000}}@{}}
\toprule\noalign{}
\begin{minipage}[b]{\linewidth}\raggedright
Class of the current state
\end{minipage} & \begin{minipage}[b]{\linewidth}\raggedright
Letters the driver may propose
\end{minipage} \\
\midrule\noalign{}
\endhead
\bottomrule\noalign{}
\endlastfoot
\texttt{Exact\allowbreak Repair} (\(\Phi_1 = 0\), \(\Gamma_2 = 0\)) & \texttt{E}
then \texttt{T}; \texttt{I} \\
\texttt{Closed\allowbreak Residual} (\(\Phi_1 > 0\), \(\Gamma_2 = 0\)) & \texttt{R},
\texttt{Λ} (Law 4), \texttt{T} to the residual terminal; \texttt{I} \\
\texttt{Active\allowbreak Closure\allowbreak Failure} (\(\Gamma_2 > 0\)) & \texttt{R},
\texttt{C}, \texttt{D}, \texttt{M}; \texttt{I}; never \texttt{T} \\
detector inadequate & \texttt{Λ}, \texttt{M}, \texttt{C}; \texttt{I} \\
\end{longtable}

The four drivers of the first draft read as follows under the laws.
\texttt{C·R·E·I·T} (autopoietic) is lawful on an exact defect.
\texttt{C·R·M·T} (refusal) is lawful when the refused state has
\(\Gamma_2 = 0\); on an open state the lawful record is \texttt{C·R·M}.
\texttt{C·R·Λ·I·T} (self-examination) is lawful from a closed survivor
with a licence when the lifted state has \(\Gamma_2 = 0\).
\texttt{D·I·T} (allopoietic) is lawful with a loss ledger when the
target has \(\Gamma_2 = 0\). The residual-progress guard of the first
draft is Theorem 11.6.

\section{13. Open laws}\label{open-laws}

These remain open. Bounded checks are recorded as computed fixtures,
never as proofs.

\begin{itemize}
\tightlist
\item
  \textbf{Trace equivalence} (DDC-OPEN-01). When do two tapes denote the
  same handling: same terminal, same diagnostic effect, same receipts,
  same loss ledger, same carrier endpoint? No equivalence is defined
  here.
\item
  \textbf{Normal forms} (DDC-OPEN-02). Whether finite tapes reduce to
  forms such as \texttt{C·R·E·I·T}, \texttt{C·R·M·T}, \texttt{C·R·Λ·I·T}
  and \texttt{D·I·T}. No universal normal form is assumed.
\item
  \textbf{Detector adequacy in general.} Demo 4 covers the constant
  detector; the empty carrier, acyclic collapse, shape-only verifiers
  and all-green or all-open terminalizers remain future fixtures, not
  results.
\item
  \textbf{Separation from rewriting logic}, \textbf{lift minimality} and
  \textbf{universal obstruction completeness} are research questions
  (Part IV).
\end{itemize}

Composition (Theorem 11.1) and the non-authority invariant (Theorem
11.5), open in the first draft, are now proved.

\section{14. Deformations of the
window}\label{deformations-of-the-window}

\textbf{Proposition 14.1} (OE010-05). A deformation
\(d_i \mapsto d_i + \mu_i\) keeps the window a complex iff
\((d_1 + \mu_1)(d_0 + \mu_0) = 0\), i.e.~iff
\(d_1 \mu_0 + \mu_1 d_0 + \mu_1 \mu_0 = 0\) (the Maurer--Cartan
condition).

\emph{Proof.} Expand and use \(d_1 d_0 = 0\). \(\square\)

\texttt{cochain.\allowbreak deformation\_obstruction} computes the left-hand side
exactly;
\texttt{test\_cochain.\allowbreak py::test\_maurer\_cartan\_obstruction\_detects\_unlawful\_deformations}
shows a zero deformation passing and a single-entry bump failing.

\section{Part IV --- Positioning}\label{part-iv-positioning}

Part IV maps the research neighbourhood. Its statements are positioning,
not verified literature conclusions; no external citation is upgraded to
a result.

\section{15. Research agenda}\label{research-agenda}

Six modules steer further work, each recast as a directed attack rather
than a validation request.

\begin{longtable}[]{@{}
  >{\raggedright\arraybackslash}p{(\linewidth - 4\tabcolsep) * \real{0.2000}}
  >{\raggedright\arraybackslash}p{(\linewidth - 4\tabcolsep) * \real{0.4000}}
  >{\raggedright\arraybackslash}p{(\linewidth - 4\tabcolsep) * \real{0.4000}}@{}}
\toprule\noalign{}
\begin{minipage}[b]{\linewidth}\raggedright
Module
\end{minipage} & \begin{minipage}[b]{\linewidth}\raggedright
Question
\end{minipage} & \begin{minipage}[b]{\linewidth}\raggedright
Relevance to DDC
\end{minipage} \\
\midrule\noalign{}
\endhead
\bottomrule\noalign{}
\endlastfoot
A. Higher categories and operads & Which structures model finite control
loops without unpriced coherence? & may supply composition laws for
tapes; not proof that DDC is an operad \\
B. Sparse \(\ell^1\) cochains & How stable are quotient obstructions
under perturbation? & diagnostic engines for \(\Phi_1\) and
\(\Gamma_2\) \\
C. Sheaf obstruction & Which Čech or cellular frameworks separate local
consistency from global coherence? & detectors for local-green,
global-red failures, once covers are declared \\
D. Verified rewriting & How can compiler passes be receipt-preserving
operations? & concrete software examples for \texttt{E}, \texttt{D},
\texttt{M}, \texttt{T} \\
E. Finite shadows of continua & How can infinite objects be admitted
only through finite shadows or oracles? & continuum claims route to
\texttt{Oracle\allowbreak Needed} \\
F. Prior art and negative results & Has anyone typed the role a
computation plays in an obstruction's life? & tests the novelty
statement \\
\end{longtable}

Every result from these modules is to be classified as an established
finite-hypothesis theorem, a plausible formalization target, an analogy,
a likely category error, or unsupported; with the exact construction,
the missing adapter data, the smallest demo and the likely
counterexample.

\section{16. Prior art and the rewriting-logic
objection}\label{prior-art-and-the-rewriting-logic-objection}

Five traditions are closest (DDC-POS-01).

\begin{itemize}
\tightlist
\item
  \textbf{Rewriting logic and Maude} is the strongest computational
  neighbour: it specifies systems by labelled, potentially conditional
  rules \(t \to t'\) modulo equations. Such a specification can include
  obstruction diagnostics, witness conditions and receipt fields. DDC
  supplies one declared combination of these choices; this paper proves
  no separation in expressive power.
\item
  \textbf{Constructor Theory} states physics through possible and
  impossible transformations; it is a philosophical analogy with no
  operation ledger or cochain quotient.
\item
  \textbf{Structural realism} aligns with the representation-first
  stance but is not mathematical support by itself.
\item
  \textbf{Hypergraph rewriting} treats evolution as rule application;
  its physical claim is empirical, not runtime verification.
\item
  \textbf{Coalgebra} gives behaviour and observation semantics, a
  promising target for trace semantics, without a cochain obstruction
  calculus.
\end{itemize}

Adjacent fields supply candidate law sources: trace monoids and abstract
rewriting (normal forms, confluence), operational semantics and effect
systems (typed letters), provenance semirings and proof-carrying code
(receipts), event structures (concurrency), double categories and string
diagrams (transport), cellular sheaves and chain homotopy (trace
equivalence), discrete Morse theory (cancellation), CEGAR and abstract
interpretation (detector adequacy), and equality saturation (repair
search).

The sharpest objection is ``DDC is rewriting logic with labels''. The
concrete contribution evaluated here is the selected policy: obstruction
diagnostics (§3), typed step guards (§7), chained records (§5), declared
licence and ledger conditions (Laws 4 and 6), refusal control (Theorem
11.3), and terminal absorption (Theorem 11.2). Theorem 11.5 proves that
emitted receipt fields say \texttt{none}; authorization remains outside
the model. These choices do not prove that rewriting logic cannot
express DDC. A formal translation and a prior-art review of the
combination remain open.

\section{17. Novelty statement}\label{novelty-statement}

Unsafe: ``DDC is a theory of everything'', ``every process is
obstruction elimination'', ``representation is all that exists''.

Defensible programme claim (DDC-POS-02): DDC formalizes a finite
operation ledger in which every structure-creating,
obstruction-exposing, obstruction-eliminating, obstruction-carrying or
obstruction-refusing event is a typed step with a replayable record that
confers no authorization. What may be distinctive is the combination of
six parts:

\begin{enumerate}
\def\labelenumi{\arabic{enumi}.}
\tightlist
\item
  finite cochain diagnostics;
\item
  typed letters recording what a step did to an obstruction;
\item
  replayable receipts;
\item
  proof-assistant and runtime counterparts;
\item
  machine-readable governance boundaries;
\item
  a firewall that stops successful representation becoming trusted
  action.
\end{enumerate}

This stays a programme claim until the prior-art review of Module F is
complete.

\section{18. Risks and fences}\label{risks-and-fences}

\begin{longtable}[]{@{}
  >{\raggedright\arraybackslash}p{(\linewidth - 2\tabcolsep) * \real{0.2727}}
  >{\raggedright\arraybackslash}p{(\linewidth - 2\tabcolsep) * \real{0.7273}}@{}}
\toprule\noalign{}
\begin{minipage}[b]{\linewidth}\raggedright
Risk
\end{minipage} & \begin{minipage}[b]{\linewidth}\raggedright
Fence
\end{minipage} \\
\midrule\noalign{}
\endhead
\bottomrule\noalign{}
\endlastfoot
Universality becomes tautology & an obstruction needs a declared
carrier, protected distinction, detector, witness or refusal rule, and
terminalizer; DDC must say when it returns \texttt{No\allowbreak Rule},
\texttt{Detector\allowbreak Inadequate} or \texttt{Oracle\allowbreak Needed} \\
``The universe is not creative'' needs a negative result & DDC tests
whether domains reuse the same letters; it does not prove that no new
structure exists \\
The word ``algebra'' overclaims & the constructed object is a typed
trace calculus with proved tape laws; equivalence and normal forms stay
open \\
Continuum smuggling & continuum objects need finite shadows, completion
receipts or oracle interfaces, else \texttt{Oracle\allowbreak Needed} \\
Authority laundering & receipts record conditional decisions;
authorization is outside this model. Theorem 11.5 establishes only the
constant receipt field \\
Detector inadequacy & a blind detector cannot seal
(\texttt{Blind\allowbreak Detector}); authority-sensitive uses need a declared
adequacy record \\
\end{longtable}

\section{19. Prior-art questions}\label{prior-art-questions}

The standing question for literature search: \emph{has any theory
defined an operation alphabet whose symbols classify the role a finite
computation played in an obstruction's life, rather than recording that
a transition occurred?} Compare against rewriting logic, trace monoids,
operational semantics, effect systems, proof-carrying code, provenance
semirings, event structures, double categories, string diagrams,
cellular sheaves, chain homotopy, deformation complexes, discrete Morse
theory, process calculi, higher rewriting, coalgebra, Constructor
Theory, hypergraph rewriting and ontic structural realism. For each,
classify and return the construction, the finite hypotheses, the missing
adapter data, the smallest demo and the likely counterexample. If DDC is
rewriting logic with labels, say exactly why; if not, name the minimal
formal object that separates them.

\section{20. Final position}\label{final-position}

DDC is a finite typed trace calculus for obstruction-handling
operations. Its product value is that it records what kind of finite
operation was performed on an obstruction, not only whether a system
passed. Its research value is a possible substrate-independent operation
ledger linking cochain obstruction, rewriting, proof-carrying traces,
provenance and authority firewalls. The warning stands: do not sell the
universe; prove trace equivalence, detector adequacy and the separation
from rewriting logic first.

\begin{Shaded}
\begin{Highlighting}[]
\NormalTok{DDC tape = replay evidence · DDC receipt = bounded verification data}
\NormalTok{DDC terminal = routing terminal · DDC authority = None}
\end{Highlighting}
\end{Shaded}

\section{21. Residuals}\label{residuals}

\begin{longtable}[]{@{}
  >{\raggedright\arraybackslash}p{(\linewidth - 4\tabcolsep) * \real{0.1333}}
  >{\raggedright\arraybackslash}p{(\linewidth - 4\tabcolsep) * \real{0.5333}}
  >{\raggedright\arraybackslash}p{(\linewidth - 4\tabcolsep) * \real{0.3333}}@{}}
\toprule\noalign{}
\begin{minipage}[b]{\linewidth}\raggedright
ID
\end{minipage} & \begin{minipage}[b]{\linewidth}\raggedright
Residual
\end{minipage} & \begin{minipage}[b]{\linewidth}\raggedright
Route
\end{minipage} \\
\midrule\noalign{}
\endhead
\bottomrule\noalign{}
\endlastfoot
R-DDC-01 & This version's new mathematics: the tape automaton and Part
III & discharged here \\
R-DDC-02 & The historical comparison model's \texttt{E} law admits
class-crossing edits; the paper requires an exact class before
\texttt{E} & historical comparison gap; resolved in governance
contracts; Rust parity separate \\
R-DDC-03 & The historical comparison model's \texttt{R} law does not fix
\(\Gamma_2\) & historical comparison gap; resolved in governance
contracts; Rust parity separate \\
R-DDC-04 & The historical comparison model does not enforce Law 8 (no
scope laundering by \texttt{C}) & historical comparison gap; resolved in
governance contracts; Rust parity separate \\
R-DDC-05 & The historical comparison model does not check Laws 3--6, the
refusal vocabulary or detector adequacy at step level & historical
comparison gap; resolved in governance contracts; Rust parity
separate \\
R-DDC-06 & \(\Phi_1\) needed a dual witness, not only a primal potential
& discharged here: \texttt{cochain.\allowbreak phi1\_certificate} \\
R-DDC-07 & Proposition 3.2 and Theorem 3.3 have written proofs only;
mechanizing them needs linear algebra beyond core Lean & Lean: a later
version \\
\end{longtable}

\section{References}\label{references}

See \hyperref[references-and-their-role]{\texttt{docs/\allowbreak references.\allowbreak md}}.

\section{Chapter 1 --- Letter E: exact repair}\label{chapter-1}

E removes a defect that lies in the image of \(d_0\). It is the
\emph{exact} part of the Hodge split, not the co-exact part; the first
draft of this chapter had the two reversed.

\subsection{1.1 Definition}\label{definition-1}

Let \(K = (V, E, F)\) be a finite oriented 2-complex with

\[(d_0 \varphi)(u \to v) = \varphi(v) - \varphi(u), \qquad
(d_1 x)(f) = \sum_{e \in \partial f} [f : e]\, x(e),\]

where \([f : e] \in \{-1, 1\}\) is the orientation of \(e\) in the
boundary of \(f\). A defect \(x \in C^1(K; \mathbb{Q})\) is
\emph{exact-repairable} when \(x \in \operatorname{im} d_0\),
equivalently \(\Phi_1(x) = 0\).

\textbf{Proposition 1.1} (DDC-E-01). An exact defect has a potential
\(\varphi\) with \(d_0 \varphi = x\), unique up to adding one constant
on each connected component. \texttt{Exact\allowbreak Repair\allowbreak Solver} finds it by
breadth-first integration from one gauge vertex per component.

\emph{Proof.} Two potentials of the same defect differ by \(\psi\) with
\(d_0 \psi = 0\), that is, \(\psi\) is constant along every edge and
hence on every component. \(\square\)

The repair is \(x \mapsto x - d_0 \varphi = 0\).

\subsection{1.2 Step law}\label{step-law}

E keeps the carrier, starts from an exact class, and leaves
\(\Phi_1 = \Gamma_2 = 0\) and a zero defect with an adequate detector.
The historical comparison model's \texttt{verify\_step} checks the
carrier and the exact class after the step
(\texttt{Exact\allowbreak Repair\allowbreak Left\allowbreak Residual}); the class before the step, the zero
defect and the detector are checked only by this package (residual
R-DDC-02).

\subsection{1.3 Receipt}\label{receipt}

\texttt{Exact\allowbreak Repair\allowbreak Receipt} records \texttt{success},
\texttt{input\_norm}, \texttt{residual\_norm}, \texttt{gauge\_vertex},
\texttt{potential}, the repaired defect (\texttt{repaired}, all zeros on
success), \texttt{error} and \texttt{authority\_effect\ =\ "None"}.

\subsection{1.4 Results}\label{results}

\textbf{Theorem 1.2} (DDC-E-02). If \(d_0 \varphi = x\) then
\(x - d_0 \varphi = 0\). Lean:
\texttt{Ddc.\allowbreak E\_gauge\_repair\_discharges}.

\textbf{Theorem 1.3} (DDC-E-03). A gauge move
\(x \mapsto x - d_0 \varphi\) changes neither \(d_1 x\) nor whether
\(x\) is exact. Consequently no gauge repair can cross a class: a defect
that is closed but not exact, or open, stays so. Lean:
\texttt{Ddc.\allowbreak E\_gauge\_preserves\_closure},
\texttt{Ddc.\allowbreak E\_gauge\_preserves\_exactness}, over an abstract complex
with additive \(d_0, d_1\) and \(d_1 d_0 = 0\).

\textbf{Theorem 1.4} (DDC-E-04). A defect with non-zero circulation
around any cycle of the declared 1-skeleton is not in
\(\operatorname{im} d_0\), and E refuses it with
\texttt{REFUSED\_NON\_ZERO\_CURL}. The name is kept for wire
compatibility; its meaning is exactly ``not in
\(\operatorname{im} d_0\)''. Lean example:
\texttt{Ddc.\allowbreak E\_triangle\_open\_not\_exact} (every edge of a filled
triangle carries 1). Tests:
\texttt{test\_exact\_repair.\allowbreak py::test\_exact\_repair\_refusal\_non\_zero\_curl},
\texttt{test\_letter\_modules.\allowbreak py::test\_exact\_repair\_refuses\_any\_cycle\_with\_circulation}.

\textbf{Theorem 1.5} (DDC-E-05). Exact repair is idempotent: repairing
the repaired (zero) defect again changes nothing. Lean:
\texttt{Ddc.\allowbreak E\_repair\_idempotent}.

The orthogonality of the exact part to the harmonic part is Theorem 3.3
of the spine.

\subsection{1.5 Falsifiers}\label{falsifiers}

\begin{itemize}
\tightlist
\item
  Demo 1 applies \texttt{E} to an open defect; the step is refused
  because the class before \texttt{E} is not exact
  (\texttt{Exact\allowbreak Repair\allowbreak From\allowbreak Non\allowbreak Exact\allowbreak Class}).
\item
  A class-crossing edit (an open state rewritten to zero) passes the
  historical comparison model but not the paper
  (\texttt{test\_automaton.\allowbreak py::test\_engine\_checks\_less\_than\_the\_paper}).
\end{itemize}

\subsection{1.6 Nonclaims}\label{nonclaims}

\begin{Shaded}
\begin{Highlighting}[]
\NormalTok{E does not claim semantic truth or the correctness of external code or models.}
\NormalTok{E is a local algebraic gauge adjustment on a declared carrier.}
\NormalTok{No authorization is conferred.}
\end{Highlighting}
\end{Shaded}

\section{Chapter 2 --- Letter R: residual projection}\label{chapter-2}

Not every defect can or should be removed. R records what remains,
keeping its class, so that an unrepaired residual is never silently
discarded or averaged away.

\subsection{2.1 Definition}\label{definition-2}

Over \(\mathbb{Q}\) with the standard inner product the 1-cochains split
as
\(C^1 = \operatorname{im} d_0 \oplus \operatorname{im} d_1^{\top} \oplus
\mathcal{H}^1\) (spine, Theorem 3.3).

\textbf{Proposition 2.1} (DDC-R-01). Every defect splits uniquely into
exact, co-exact and harmonic parts. The first draft named only two
parts, which is correct only when there are no faces. Evidence: spine
Theorem 3.3, \texttt{cochain.\allowbreak hodge\_split},
\texttt{test\_cochain.\allowbreak py::test\_hodge\_split\_is\_orthogonal\_and\_matches\_the\_classes}.

R's canonical representative subtracts the coboundary of the potential
that matches the defect on a spanning forest. What remains is supported
on the edges outside the forest: one value per independent cycle.

\textbf{Proposition 2.2} (DDC-R-06). The register has
\(b_1 = |E| - |V| + c\) coordinates for a graph with \(c\) components.
They are the circulations of the defect around the fundamental cycles,
and all of them vanish exactly when the defect is exact. The first draft
handled a single cycle only. \texttt{Graph\allowbreak Residual\allowbreak Projector} covers any
finite graph, and \texttt{Residual\allowbreak Projector} keeps the single-cycle
case. Tests:
\texttt{test\_letter\_modules.\allowbreak py::test\_graph\_projector\_keeps\_one\_coordinate\_per\_independent\_cycle},
\texttt{test\_cochain.\allowbreak py::test\_a\_hollow\_cycle\_carries\_a\_closed\_residual}.

\subsection{2.2 Step law}\label{step-law-1}

R keeps the carrier, \(\Phi_1\) and \(\Gamma_2\), with an adequate
detector. The historical comparison model checks the carrier and
\(\Phi_1\) (\texttt{Residual\allowbreak Projection\allowbreak Changed\allowbreak Mass}); the \(\Gamma_2\)
clause is paper-only (residual R-DDC-03).

\subsection{2.3 Receipt}\label{receipt-1}

\texttt{Residual\allowbreak Projection\allowbreak Receipt} records the raw, exact and residual
norms, both components, the status (\texttt{RESIDUAL\_RETAINED} or
\texttt{EXACT\_TENSION\_ONLY}), the cycle coordinates and
\texttt{authority\_effect\ =\ "None"}.

These norms are unweighted \(\ell^1\) norms of the recorded vectors,
rather than certificates of \(\Phi_1\). A fixed forest representative
may have norm strictly greater than the quotient minimum. For example,
on the increasingly oriented complete graph with vertices \(a<b<c<d\)
and value \(-1\) on all six edges, the forest rooted at \(a\) retains
\((-1,-1,-1)\) on \(bc,bd,cd\). Its norm is 3; the potential
\((0,-1,-1,-2)\) attains objective 2 and a dual witness certifies
\(\Phi_1=2\). Test:
\texttt{test\_publication\_regressions.\allowbreak py::test\_forest\_register\_norm\_is\_not\_the\_quotient\_minimum}.
R preserves the quotient distance because its difference is exact,
regardless of whether its representative attains the minimum.

\texttt{evaluate\_xi\_tie\_break} compares a secondary exact score only
after equal primary diagnostics. It checks the supplied diagnostic
values and makes no claim to authenticate or remeasure the candidates.

\subsection{2.4 Results}\label{results-1}

\textbf{Theorem 2.3} (DDC-R-02). The canonical R move is a gauge move,
so it keeps \(\Gamma_2\) and exactness; a lawful R step keeps the class.
Lean: \texttt{Ddc.\allowbreak R\_gauge\_move\_is\_lawful},
\texttt{Ddc.\allowbreak R\_keeps\_diagnostics}, \texttt{Ddc.\allowbreak R\_keeps\_class}.

\textbf{Proposition 2.4} (DDC-R-03). The defect is conserved edge by
edge: the exact component plus the residual component equals the raw
defect. Test:
\texttt{test\_letter\_modules.\allowbreak py::test\_projection\_conserves\_the\_defect}.

\textbf{Proposition 2.5} (DDC-R-04). Projecting the representative again
returns it unchanged. Test:
\texttt{test\_letter\_modules.\allowbreak py::test\_graph\_projector\_keeps\_one\_coordinate\_per\_independent\_cycle}.

\textbf{Proposition 2.6} (DDC-R-05). Terminalizing an unchanged
non-exact state never yields \texttt{Check\allowbreak Passed}. A closed residual is
sealed only on the residual terminal, and an open one not at all (spine,
Theorem 11.7). Test:
\texttt{test\_automaton.\allowbreak py::test\_run\_tape\_stops\_at\_the\_first\_unlawful\_step\_and\_chains\_receipts}.

\subsection{2.5 Falsifiers}\label{falsifiers-1}

\begin{itemize}
\tightlist
\item
  An R step that lowers \(\Gamma_2\) is refused
  (\texttt{Residual\allowbreak Projection\allowbreak Changed\allowbreak Closure}).
\item
  An R step that moves \(\Phi_1\) is refused by the paper and the
  historical comparison model (\texttt{Residual\allowbreak Projection\allowbreak Changed\allowbreak Mass}).
\end{itemize}

The abstract step law compares diagnostics only. Equal norms do not
identify quotient classes: on a hollow triangle, (1,0,0) and (−1,0,0)
have equal diagnostics but different circulation.
\texttt{Graph\allowbreak Residual\allowbreak Projector} subtracts an actual coboundary;
\texttt{run\_tape(.\allowbreak .\allowbreak .\allowbreak ,\ carriers=.\allowbreak .\allowbreak .\allowbreak )} checks that stronger condition. A
booked descent can subsequently change a non-exact state into an exact
one.

\subsection{2.6 Nonclaims}\label{nonclaims-1}

\begin{Shaded}
\begin{Highlighting}[]
\NormalTok{R does not resolve the defect and does not claim that a retained residual will vanish.}
\NormalTok{R binds unresolved residual mass into the tape. No authorization is conferred.}
\end{Highlighting}
\end{Shaded}

\section{Chapter 3 --- Letter C: context exposure}\label{chapter-3}

Evaluations are local. C restricts the ambient carrier to a window and
exposes every edge that straddles it, so that external dependencies
become declared boundary conditions instead of hidden entanglement.

\subsection{3.1 Definition}\label{definition-3}

For a finite graph \(K = (V, E)\) and a window \(V_W \subseteq V\), the
internal edges are those with both ends in \(V_W\). The boundary
interface \(\partial_E W\) is the set of edges with exactly one end in
\(V_W\). C maps a cochain \(c\) to \((c|_W, c|_{\partial_E W})\).

\subsection{3.2 Step law}\label{step-law-2}

C changes the carrier and never lowers \(\Phi_1\) or \(\Gamma_2\) (Law
8). The historical comparison model checks the carrier change
(\texttt{Carrier\allowbreak Unchanged\allowbreak Under\allowbreak Restructure}); the monotonicity clause is
paper-only (residual R-DDC-04). Hiding an obstruction by choosing a
smaller window is scope laundering; it has to be priced as a descent
with a loss ledger or a licensed lift.

\subsection{3.3 Receipt}\label{receipt-2}

\texttt{Context\allowbreak Exposure\allowbreak Receipt} records the window, the internal and
boundary edges, both restricted cochains, the status
\texttt{CONTEXT\_EXPOSED} and \texttt{authority\_effect\ =\ "None"}.

\subsection{3.4 Results}\label{results-2}

\textbf{Proposition 3.1} (DDC-C-01). The internal edges, the boundary
edges and the remaining edges partition \(E\). Test:
\texttt{test\_context\_exposure.\allowbreak py::test\_context\_window\_disjoint\_partition}.

\textbf{Proposition 3.2} (DDC-C-02). An edit supported on the internal
edges leaves the boundary cochain unchanged. Test:
\texttt{test\_letter\_modules.\allowbreak py::test\_exposure\_is\_read\_only\_and\_internal\_edits\_keep\_the\_boundary}.

\textbf{Proposition 3.3} (DDC-C-03). Two exposures that agree on every
shared edge glue to a unique cochain on the union of their exposed
edges; \texttt{glue} refuses when a shared edge disagrees. Test:
\texttt{test\_letter\_modules.\allowbreak py::test\_context\_exposure\_glues\_agreeing\_windows\_uniquely}.

\textbf{Proposition 3.4} (DDC-C-04). Exposure is read-only: it returns
restrictions and never alters the ambient cochain. Test as for DDC-C-02.

\textbf{Theorem 3.5} (DDC-C-05). The boundary of an intersection of
windows lies in the union of their boundaries. Lean:
\texttt{Ddc.\allowbreak C\_boundary\_of\_intersection}; test:
\texttt{test\_letter\_modules.\allowbreak py::test\_boundary\_of\_an\_intersection\_lies\_in\_the\_union\_of\_boundaries}.

\textbf{Theorem 3.6} (DDC-C-06). A lawful C step changes the carrier and
never lowers \(\Phi_1\) or \(\Gamma_2\). Lean:
\texttt{Ddc.\allowbreak C\_is\_not\_repair}, \texttt{Ddc.\allowbreak C\_changes\_carrier}.

\subsection{3.5 Falsifiers}\label{falsifiers-2}

\begin{itemize}
\tightlist
\item
  A C step that lowers \(\Gamma_2\) is refused
  (\texttt{Context\allowbreak Exposure\allowbreak Lowered\allowbreak Obstruction}).
\item
  A C step that leaves the carrier unchanged is refused by paper and
  historical comparison model.
\end{itemize}

\subsection{3.6 Nonclaims}\label{nonclaims-2}

\begin{Shaded}
\begin{Highlighting}[]
\NormalTok{C does not claim that a window is complete or that unexposed values are trivial.}
\NormalTok{C is not repair. No authorization is conferred.}
\end{Highlighting}
\end{Shaded}

\section{Chapter 4 --- Letter I: integration}\label{chapter-4}

No progress counts until it is committed. I appends a receipt to an
append-only, hash-chained tape and changes nothing else.

\subsection{4.1 Definition}\label{definition-4}

The tape at step \(t\) is \(T_t = [r_0, \dots, r_{t-1}]\). With
\(h_0 = \mathrm{SHA256}(\text{genesis root})\), the receipt for letter
\(\alpha \in \{E, R, C, I, \Lambda, D, M, T\}\) is chained as

\[h_{t+1} = \mathrm{SHA256}\big(t : h_t : \alpha : \mathrm{json}(r_t)\big).\]

Every letter enters the chain, I included; the first draft's letter set
left I out. The index, the previous head (64 hex digits) and the letter
cannot contain the separator, so the preimage parses uniquely. Payloads
are canonical JSON with exact rationals written as strings, and floats
are refused.

\subsection{4.2 Step law}\label{step-law-3}

I leaves the state unchanged and carries a receipt (Law 5). The
historical comparison model checks the unchanged state
(\texttt{Inert\allowbreak Op\allowbreak Mutated\allowbreak State}); the receipt clause is paper-only
(residual R-DDC-05).

\subsection{4.3 Receipt}\label{receipt-3}

\texttt{Tape\allowbreak Entry} records \texttt{index}, \texttt{prev\_hash},
\texttt{letter}, the canonical \texttt{payload} and
\texttt{current\_hash}.

\subsection{4.4 Results}\label{results-3}

\textbf{Theorem 4.1} (DDC-I-01). The tape only grows: a tape is a prefix
of every later tape, and the receipts of a replay are a prefix of the
receipts of any extension of it. Lean: \texttt{Ddc.\allowbreak I\_commit\_prefix},
\texttt{Ddc.\allowbreak I\_receipts\_prefix}.

\textbf{Proposition 4.2} (DDC-I-02). A broken hash or predecessor link
is located by \texttt{first\_broken\_index}. Replacing a final entry and
rehashing it, rewriting a suffix, or truncating the tape can leave a
consistent chain. \texttt{verify\_integrity(checkpoint)} compares with
an independently retained genesis hash, head hash and length; under
SHA-256 collision resistance it detects these changes. A checkpoint
obtained from the suspect tape supplies no authentication. A final
checkpoint detects a mismatch but does not locate the original first
alteration. Tests:
\texttt{test\_review\_regressions.\allowbreak py::test\_tape\_checkpoint\_detects\_rehashed\_replacement\_and\_truncation}.

\textbf{Proposition 4.3} (DDC-I-03). Replaying the same steps from the
same state reproduces every state and the same head hash. Tests:
\texttt{test\_automaton.\allowbreak py::test\_replay\_is\_deterministic}, and the
pinned heads in \texttt{test\_fsm\_and\_demos.\allowbreak py}.

\textbf{Proposition 4.4} (DDC-I-04). A step that fails its law is not
committed: replay stops before recording it. Lean:
\texttt{Ddc.\allowbreak I\_needs\_receipt}; test:
\texttt{test\_automaton.\allowbreak py::test\_run\_tape\_stops\_at\_the\_first\_unlawful\_step\_and\_chains\_receipts}.

\subsection{4.5 Falsifiers}\label{falsifiers-3}

\begin{itemize}
\tightlist
\item
  An I step without a receipt is refused (\texttt{Missing\allowbreak Receipt}).
\item
  An I step that alters the state is refused by paper and historical
  comparison model.
\end{itemize}

\subsection{4.6 Nonclaims}\label{nonclaims-3}

\begin{Shaded}
\begin{Highlighting}[]
\NormalTok{I does not claim semantic authority or that committed steps are optimal.}
\NormalTok{I provides deterministic commitment; historical integrity needs a trusted checkpoint.}
\NormalTok{No authorization is conferred.}
\end{Highlighting}
\end{Shaded}

\section{Chapter 5 --- Letter Λ: structural lift}\label{chapter-5}

When an obstruction cannot be handled on the current carrier, ad hoc
expansion would launder authority. Λ is licensed carrier expansion: it
needs a licence, and it needs a trigger that makes a lift the right
move. A trigger is not a licence.

\subsection{5.1 Definition}\label{definition-5}

The lift pads with zeros: \(\iota : \mathbb{Q}^n \hookrightarrow
\mathbb{Q}^{n+k}\), \(\iota(x)_i = x_i\) for \(i \le n\) and \(0\)
otherwise.

\subsection{5.2 Step law}\label{step-law-4}

Λ changes the carrier, carries a licence, and starts from a
\texttt{Closed\allowbreak Residual} or from an inadequate detector (Law 4). The
historical comparison model checks the carrier change; the licence and
the trigger are paper-only (residual R-DDC-05). The first draft did not
state the trigger.

\subsection{5.3 Receipt}\label{receipt-4}

\texttt{Lift\allowbreak Receipt} records the initial and extended dimensions, the
licence, the lifted vector, the trigger, \texttt{success},
\texttt{error} (\texttt{REFUSED\_UNLICENSED\_LIFT},
\texttt{REFUSED\_LIFT\_TRIGGER} or
\texttt{INVALID\_NEGATIVE\_EXPANSION}) and
\texttt{authority\_effect\ =\ "None"}.

The auxiliary \texttt{Finite\allowbreak Lambda\allowbreak Ledger} stores exact nonnegative
diagnostic rows and checks endpoint matching, adjacency, non-final
\texttt{is\_descent} flags and the final \texttt{is\_valid} flag. Both
flags must be actual booleans; the row sequence is snapshotted. They
remain assumptions, rather than checked preservation or licence
certificates. This consistency gate does not execute the external Rust
\texttt{require\_descent\_step} or \texttt{require\_valid} methods and
establishes no current contract conformance. It also does not impose a
resolved final endpoint on a nonempty ledger; that condition needs a
separate policy.

\subsection{5.4 Results}\label{results-4}

\textbf{Theorem 5.1} (DDC-L-01). Zero-padding preserves the \(\ell^1\)
norm, and every inner product among old vectors. Lean:
\texttt{Ddc.\allowbreak L\_pad\_preserves\_l1}; test:
\texttt{test\_letter\_modules.\allowbreak py::test\_lift\_trigger\_law\_and\_norm\_preservation}.

\textbf{Proposition 5.2} (DDC-L-02). A lawful Λ carries a licence and a
trigger. Lean: \texttt{Ddc.\allowbreak L\_needs\_licence\_and\_trigger}.

\textbf{Proposition 5.3} (DDC-L-03). The new coordinates are orthogonal
to the image of \(\iota\), since their supports are disjoint.

\textbf{Proposition 5.4} (DDC-L-04). An unlicensed or untriggered lift
fails closed and returns the original vector. Tests:
\texttt{test\_structural\_lift.\allowbreak py::test\_unlicensed\_lift\_refusal},
\texttt{test\_letter\_modules.\allowbreak py::test\_lift\_trigger\_law\_and\_norm\_preservation}.

\textbf{Theorem 5.5} (DDC-TAPE-09). Let \(K_0\) sit in \(K_1\) with
restriction maps that commute with \(d_0\), and let \(x'\) on \(K_1\)
restrict to \(x\) on \(K_0\). If \(x'\) is exact on \(K_1\), then \(x\)
is exact on \(K_0\). So a non-exact residual stays non-exact under every
lift that extends it. Lean: \texttt{Ddc.\allowbreak L\_cannot\_exactify}.

\emph{Proof.} If \(x' = d_0 \varphi\) then
\(x = \rho_1(d_0 \varphi) = d_0(\rho_0 \varphi)\). \(\square\)

Demo 3 shows both halves. A lift that adds capacity away from the cycle
keeps \(\Phi_1 = 1\), and the survivor is sealed on the residual
terminal. A lift that attaches a face bounding the cycle turns the
survivor into an open defect (\(\Gamma_2 = 1\)) that \texttt{T} may not
seal.

\subsection{5.5 Falsifiers}\label{falsifiers-4}

\begin{itemize}
\tightlist
\item
  A lift without a licence is refused (\texttt{Unlicensed\allowbreak Lift}).
\item
  A lift from an exact or an open, adequately detected state is refused
  (\texttt{Lift\allowbreak Without\allowbreak Trigger}).
\end{itemize}

\subsection{5.6 Nonclaims}\label{nonclaims-4}

\begin{Shaded}
\begin{Highlighting}[]
\NormalTok{Λ does not grant unbounded or continuous expansion and does not resolve a survivor.}
\NormalTok{Λ requires an explicit discrete licence. No authorization is conferred.}
\end{Highlighting}
\end{Shaded}

\section{Chapter 6 --- Letter D: descent}\label{chapter-6}

Moving to a coarser carrier discards distinctions. D does not claim that
the coarse state is faithful; it books exactly what was lost.

\subsection{6.1 Definition}\label{definition-6}

With \(\pi : \mathbb{Q}^{n+k} \to \mathbb{Q}^n\) keeping the first \(n\)
coordinates and \(\iota\) the zero-padding inclusion, the loss vector is
\(x - \iota(\pi(x))\) (the dropped tail) and the loss scalar is its
\(\ell^1\) norm \(L(x) = \sum_{i > n} \lvert x_i \rvert\).

\subsection{6.2 Step law}\label{step-law-5}

D changes the carrier and carries a loss ledger (Law 6). The historical
comparison model checks the carrier change; the ledger clause is
paper-only (residual R-DDC-05). D is one of the two letters that may
lower \(\Phi_1\) or \(\Gamma_2\) (spine, Theorem 11.6), which is why its
loss must be booked.

\subsection{6.3 Receipt}\label{receipt-5}

\texttt{Descent\allowbreak Receipt} records the source and target dimensions, the
coarse vector, the loss vector, the loss scalar, \texttt{is\_lossless},
the status (\texttt{LOSSLESS\_DESCENT} or \texttt{DESCENT\_WITH\_LOSS})
and \texttt{authority\_effect\ =\ "None"}. \texttt{ledger\_entry} is the
ledger a D step carries.

\subsection{6.4 Results}\label{results-5}

\textbf{Proposition 6.1} (DDC-D-01). The loss is non-negative (a sum of
absolute values; in Lean a natural number).

\textbf{Theorem 6.2} (DDC-D-02). The descent is lossless exactly when
every dropped coordinate is zero, that is, when
\(x \in \iota(\mathbb{Q}^n)\). Lean: \texttt{Ddc.\allowbreak D\_lossless\_iff},
\texttt{Ddc.\allowbreak l1\_eq\_zero\_iff}.

\textbf{Theorem 6.3} (DDC-D-03). Projecting twice is projecting once.
Lean: \texttt{Ddc.\allowbreak D\_projection\_idempotent}.

\textbf{Proposition 6.4} (DDC-D-04). A non-zero loss can never produce a
\texttt{LOSSLESS\_DESCENT} receipt, and a lawful D carries its ledger.
Lean: \texttt{Ddc.\allowbreak D\_needs\_ledger}; tests:
\texttt{test\_descent.\allowbreak py::test\_descent\_with\_positive\_loss},
\texttt{test\_letter\_modules.\allowbreak py::test\_descent\_books\_its\_ledger}.

\subsection{6.5 Falsifiers}\label{falsifiers-5}

\begin{itemize}
\tightlist
\item
  Demo 5 without a ledger: \texttt{D} is refused
  (\texttt{Missing\allowbreak Loss\allowbreak Ledger}), and the lawful record is \texttt{M·T}
  with \texttt{REFUSED\_MISSING\_LOSS\_LEDGER}.
\item
  CREDIT without a ledger at D is not lawful (spine, Theorem 11.8).
\end{itemize}

\subsection{6.6 Nonclaims}\label{nonclaims-5}

\begin{Shaded}
\begin{Highlighting}[]
\NormalTok{D does not claim lossless projection unless the loss is zero.}
\NormalTok{D does not permit unrecorded loss of distinctions. No authorization is conferred.}
\end{Highlighting}
\end{Shaded}

\section{Chapter 7 --- Letter M: typed refusal}\label{chapter-7}

The ability to refuse is co-equal with the ability to deform. M is a
lawful refusal outcome: it leaves the state unchanged, names its reason
and moves control to \texttt{Refused}. It does not seal the tape; a
later C, Λ or T still needs its own law and control check.

\subsection{7.1 Definition}\label{definition-7}

A request is a letter, a locus, an alleged witness and a cost. Grounding
is checked at intake, before admission; witness and budget are checked
on the tape as M.

\subsection{7.2 Step law}\label{step-law-6}

M leaves the state unchanged and carries a reason from the closed
vocabulary below. The historical comparison model checks the unchanged
state (\texttt{Inert\allowbreak Op\allowbreak Mutated\allowbreak State}). It carries its reason as a free
string (\texttt{dda\_alphabet::Step\allowbreak Outcome::Refused\allowbreak To\allowbreak Execute}), so this
vocabulary is owned by the paper (residual R-DDC-05). After M, only
\texttt{C}, \texttt{Λ} or \texttt{T} may follow (spine, Theorem 11.3).

\subsection{7.3 The refusal vocabulary}\label{the-refusal-vocabulary}

Intake refusals are recorded before admission and never enter the tape
(\texttt{enters\_ddc\_trace} is false):

\begin{longtable}[]{@{}
  >{\raggedright\arraybackslash}p{(\linewidth - 2\tabcolsep) * \real{0.5000}}
  >{\raggedright\arraybackslash}p{(\linewidth - 2\tabcolsep) * \real{0.5000}}@{}}
\toprule\noalign{}
\begin{minipage}[b]{\linewidth}\raggedright
Code
\end{minipage} & \begin{minipage}[b]{\linewidth}\raggedright
Meaning
\end{minipage} \\
\midrule\noalign{}
\endhead
\bottomrule\noalign{}
\endlastfoot
\texttt{REFUSED\_UNTYPED} & the request or its defect is not typed on
the declared carrier \\
\texttt{REFUSED\_UNGROUNDED\_LOCUS} & the locus is not in the declared
carrier \\
\texttt{ADAPTER\_INADEQUATE} & the adapter cannot represent the
request \\
\end{longtable}

In-trace reasons carried by M:

\begin{longtable}[]{@{}
  >{\raggedright\arraybackslash}p{(\linewidth - 2\tabcolsep) * \real{0.5000}}
  >{\raggedright\arraybackslash}p{(\linewidth - 2\tabcolsep) * \real{0.5000}}@{}}
\toprule\noalign{}
\begin{minipage}[b]{\linewidth}\raggedright
Code
\end{minipage} & \begin{minipage}[b]{\linewidth}\raggedright
Raised when
\end{minipage} \\
\midrule\noalign{}
\endhead
\bottomrule\noalign{}
\endlastfoot
\texttt{REFUSED\_INVALID\_WITNESS} & the witness is missing or does not
verify \\
\texttt{REFUSED\_BUDGET\_EXHAUSTED} & the cost exceeds the remaining
budget \\
\texttt{REFUSED\_NON\_ZERO\_CURL} & E was requested for a defect not in
\(\operatorname{im} d_0\) \\
\texttt{REFUSED\_UNLICENSED\_LIFT} & Λ was requested without a
licence \\
\texttt{REFUSED\_LIFT\_TRIGGER} & Λ was requested without a Law 4
trigger \\
\texttt{REFUSED\_MISSING\_LOSS\_LEDGER} & D was requested without a loss
ledger \\
\texttt{REFUSED\_MISSING\_RECEIPT} & I was requested without a
receipt \\
\texttt{REFUSED\_REQUIRES\_LIFT} & a closed survivor needs a lift that
is not licensed \\
\texttt{REFUSED\_DETECTOR\_INADEQUATE} & the detector cannot certify the
diagnostics \\
\texttt{REFUSED\_OPEN\_OBSTRUCTION} & an open obstruction
(\(\Gamma_2 > 0\)) cannot be sealed \\
\end{longtable}

The first draft used both \texttt{Missing\allowbreak Loss\allowbreak Ledger} and SCREAMING\_CASE
codes, and listed the ungrounded locus both as an intake refusal and as
an M reason. It is an intake refusal.

\subsection{7.4 Results}\label{results-6}

\textbf{Theorem 7.1} (DDC-M-01). A lawful M leaves the state unchanged
and names a reason, and it moves the control to \texttt{Refused}. Lean:
\texttt{Ddc.\allowbreak M\_preserves\_state}, \texttt{Ddc.\allowbreak M\_refuses}.

\textbf{Proposition 7.2} (DDC-M-02). An ungrounded locus is refused at
intake and never enters the tape. Test:
\texttt{test\_letter\_modules.\allowbreak py::test\_ungrounded\_locus\_is\_an\_intake\_refusal\_that\_never\_enters\_the\_tape}.

\textbf{Proposition 7.3} (DDC-M-03). A missing or invalid witness yields
\texttt{REFUSED\_INVALID\_WITNESS}. Test:
\texttt{test\_typed\_refusal.\allowbreak py::test\_refusal\_invalid\_witness}.

\textbf{Proposition 7.4} (DDC-M-04). A cost above the remaining budget
yields \texttt{REFUSED\_BUDGET\_EXHAUSTED}. Test:
\texttt{test\_typed\_refusal.\allowbreak py::test\_refusal\_budget\_exhausted}.

\subsection{7.5 Falsifiers}\label{falsifiers-6}

\begin{itemize}
\tightlist
\item
  An M step without a reason, or with a reason outside the vocabulary,
  is refused (\texttt{Missing\allowbreak Refusal\allowbreak Reason}).
\item
  An M step that alters the state is refused by paper and historical
  comparison model.
\end{itemize}

\subsection{7.6 Nonclaims}\label{nonclaims-6}

\begin{Shaded}
\begin{Highlighting}[]
\NormalTok{M is not a system failure. It is a typed, lawful refusal that preserves state.}
\NormalTok{No authorization is conferred.}
\end{Highlighting}
\end{Shaded}

\section{Chapter 8 --- Letter T: terminalization}\label{chapter-8}

A recursive evaluation needs a halting law that prevents infinite
expansion and self-certification. T closes a tape on a routing terminal
when the declared terminal policy allows it, subject to the sealing
conditions below.

\subsection{8.1 Definition}\label{definition-8}

\texttt{Can\allowbreak Continue\_n} is the conjunction of the nine conjuncts named
by the historical comparison model
(\texttt{omega\_discrete\_deformation\_calculus::oeprmalg2::CONTINUATION\_CONJUNCTS}):

\begin{enumerate}
\def\labelenumi{\arabic{enumi}.}
\tightlist
\item
  an outstanding obligation exists;
\item
  a higher obligation witness exceeds it;
\item
  the next carrier is declared;
\item
  the adapter is sound (its loss ledger is booked);
\item
  the detector is adequate;
\item
  any lift is minimal and licensed;
\item
  the old skeleton is preserved;
\item
  budget is available;
\item
  no GTLA continuation refusal was issued.
\end{enumerate}

The first draft's three-condition form (positive budget, unresolved
residual, strict progress) is a weakening. It is conjuncts 8 and 1, plus
the progress guard of spine Theorem 11.6.

\subsection{8.2 Step law}\label{step-law-7}

T leaves the state unchanged, requires \(\Gamma_2 = 0\), and requires an
adequate detector. The historical comparison model checks the unchanged
state and \(\Gamma_2 = 0\) (\texttt{Terminalized\allowbreak Open\allowbreak Residual}); the
detector clause is paper-only (residual R-DDC-05). After T nothing is
admitted.

\subsection{8.3 Verdicts}\label{verdicts}

The abstract Lean T law does not include \texttt{Can\allowbreak Continue\_n}; its
theorem concerns only the conjunction itself.
\texttt{Terminalization\allowbreak Evaluator.\allowbreak evaluate} accepts optional
\texttt{continuation\_facts}. After the open/blind-state guards, if all
nine hold it returns \texttt{CAN\_CONTINUE}; if one fails, an open or
blind state still cannot seal, and a closed positive residual receives
\texttt{HALTED\_CONTINUATION\_POLICY}. These booleans are supplied
assumptions, not independently verified facts.

Without \texttt{continuation\_facts}, the separate
residual/budget/lookback policy returns one of the following. Its
simple-policy tests do not establish the nine-fact policy for a deployed
system:

\begin{longtable}[]{@{}
  >{\raggedright\arraybackslash}p{(\linewidth - 4\tabcolsep) * \real{0.3333}}
  >{\raggedright\arraybackslash}p{(\linewidth - 4\tabcolsep) * \real{0.3333}}
  >{\raggedright\arraybackslash}p{(\linewidth - 4\tabcolsep) * \real{0.3333}}@{}}
\toprule\noalign{}
\begin{minipage}[b]{\linewidth}\raggedright
Verdict
\end{minipage} & \begin{minipage}[b]{\linewidth}\raggedright
When
\end{minipage} & \begin{minipage}[b]{\linewidth}\raggedright
Routing
\end{minipage} \\
\midrule\noalign{}
\endhead
\bottomrule\noalign{}
\endlastfoot
\texttt{SEAL\_REFUSED\_OPEN\_OBSTRUCTION} & \(\Gamma_2 > 0\) & not
terminal; the tape ends in M \\
\texttt{SEAL\_REFUSED\_DETECTOR\_INADEQUATE} & blind detector & not
terminal; the tape ends in M \\
\texttt{LOCAL\_GREEN\_FOR\_SCOPE} & zero residual & \texttt{Check\allowbreak Passed}
inside the declared scope \\
\texttt{HALTED\_OUT\_OF\_BUDGET} & closed residual, budget spent &
residual terminal \\
\texttt{HALTED\_STAGNANT\_PROGRESS} & no progress over the lookback
window & residual terminal \\
\end{longtable}

Every receipt records
\texttt{authority\_terminal\ =\ "F\_term:REFUSED"}. A local green
licenses nothing: refusal on the authorization axis is an explicit
package policy, not an inferred external theorem. The first draft
emitted a terminal verdict for an out-of-budget open obstruction; that
verdict now routes to M (\texttt{REFUSED\_OPEN\_OBSTRUCTION}).

\subsection{8.4 Results}\label{results-7}

\textbf{Proposition 8.1} (DDC-T-01). If every charged step costs at
least \(\varepsilon > 0\) from an initial budget \(B_0\), at most
\(\lfloor B_0 / \varepsilon \rfloor\) charged steps occur. This is a
conditional theorem with a one-line proof: the budget is non-negative
and falls by \(\varepsilon\) per step.

\textbf{Proposition 8.2} (DDC-T-02). Under the simple policy, a zero
residual with \(\Gamma_2 = 0\) and an adequate detector yields
\texttt{LOCAL\_GREEN\_FOR\_SCOPE} and records the unlicensed authority
terminal. Tests:
\texttt{test\_terminalization.\allowbreak py::test\_terminal\_local\_green},
\texttt{test\_letter\_modules.\allowbreak py::test\_terminal\_evaluator\_refuses\_to\_seal\_open\_or\_blind\_states}.

\textbf{Proposition 8.3} (DDC-T-03). A non-decreasing residual over the
lookback window halts with \texttt{HALTED\_STAGNANT\_PROGRESS}. Test:
\texttt{test\_terminalization.\allowbreak py::test\_terminal\_stagnant\_progress}.

\textbf{Theorem 8.4} (DDC-T-04). T seals: the control becomes
\texttt{Terminalized} and admits nothing. Lean: \texttt{Ddc.\allowbreak T\_seals},
\texttt{Ddc.\allowbreak terminal\_absorbing}.

\textbf{Theorem 8.5} (DDC-T-05). A lawful T leaves the state unchanged,
has \(\Gamma_2 = 0\) and an adequate detector, and lands on the exact or
the closed class. Lean: \texttt{Ddc.\allowbreak T\_never\_seals\_open},
\texttt{Ddc.\allowbreak T\_routing}.

\textbf{Theorem 8.6} (DDC-T-06). \texttt{Can\allowbreak Continue\_n} holds exactly
when all nine conjuncts hold, and each conjunct is a strict boolean,
never an integer. Lean: \texttt{Ddc.\allowbreak can\allowbreak Continue\_iff}; test:
\texttt{test\_letter\_modules.\allowbreak py::test\_can\_continue\_is\_the\_conjunction\_of\_the\_nine\_named\_conjuncts}.

\subsection{8.5 Falsifiers}\label{falsifiers-7}

\begin{itemize}
\tightlist
\item
  Demo 2 appends \texttt{T} to an open refusal: refused
  (\texttt{Terminalized\allowbreak Open\allowbreak Residual}).
\item
  Demo 4 appends \texttt{T} under a blind detector: refused
  (\texttt{Blind\allowbreak Detector}).
\end{itemize}

\subsection{8.6 Nonclaims}\label{nonclaims-7}

T does not mint global authority and never emits an unscoped green
verdict. A terminal receipt records a conditional terminal-policy
decision. No authorization is conferred.

\newpage

\section{Appendices --- Discrete Deformation Calculus
v1.0.3}\label{appendices}

\textbf{Part of:} \hyperref[discrete-deformation-calculus]{Discrete
Deformation Calculus v1.0.3}

The appendices hold candidate adapters, observer notes and worked
trajectories. Nothing here is part of the core calculus unless its
ledger row says so.

\subsection{A. Candidate adapter: ZX diagrams to
cochains}\label{a.-candidate-adapter-zx-diagrams-to-cochains}

\begin{quote}
\textbf{Fence.} A candidate adapter from planar ZX diagrams to finite
cochains. It is not core DDC and asserts no quantum or physical result
(DDC-OPEN-03).
\end{quote}

A finite ZX diagram gives a graph carrier: spiders with a type \(Z\) or
\(X\) and a phase \(\theta_v\) are vertices, oriented wires with a phase
discrepancy \(\omega_e\) are edges, and closed plaquettes are faces. The
candidate coboundary is colour-relative:

\[(d_0 \theta)_{u \to v} = \theta_v - \theta_u - H_{u,v}, \qquad
H_{u,v} = \begin{cases} 0 & \text{same type,} \\ \pi/2 & \text{different type,} \end{cases}\]

with \(d_1\) summing signed discrepancies around each plaquette. Two
obligations are open. First, \(d_0\) with a non-zero \(H\) is affine,
not linear, so \(d_1 \circ d_0 = 0\) must be proved or refused for the
declared \(H\). Second, phases are real, so an exact treatment needs a
finite shadow over \(\mathbb{Q}\) or a declared oracle. Spider fusion
\(Z(\alpha) \circ Z(\beta) \Rightarrow
Z(\alpha + \beta)\) would read as E when a local \(\theta\) with
\(\theta_v - \theta_u = \omega_e\) exists, giving the candidate tape
\texttt{C·R·E·I·T}.

\subsection{B. Observer gauges and dual
witnesses}\label{b.-observer-gauges-and-dual-witnesses}

\begin{quote}
\textbf{Fence.} An observer-gauge note. It is not core DDC, and it
leaves \(\mathbb{Q}\) where it takes a square root (DDC-OPEN-04).
\end{quote}

For a two-component defect \(\omega = (1, 1)\) on one edge of a hollow
six-cycle, a coordinate-separable observer reads
\(\lVert (1,1) \rVert_1 = 2\), and a rotation-invariant observer reads
\(\lVert (1,1) \rVert_2 = \sqrt{2}\). The ratio \(\sqrt{2}\) is specific
to this defect and these gauges; for (1,0) the ratio is 1. A feasible
dual flow must have zero divergence. For the separable gauge, the vector
(1,1) circulates on every consistently oriented cycle edge; only its
pairing with the one-edge defect is concentrated on that edge. The
rotation-invariant gauge similarly uses the constant unit flow
\((1,1)/\sqrt{2}\) on every edge. A one-edge dual vector would violate
flow conservation. Formalizing the finite dual certificate, with its
Fenchel--Rockafellar hypotheses named, is open.

\subsection{C. Worked trajectories}\label{c.-worked-trajectories}

\textbf{C.1 Trajectory recording} (OE010-06). A trajectory maps concrete
tool and human actions into cochain classes, and cochains back into
readable evidence, with \(d_{k} \circ d_{k-1} = 0\) on the window. This
is a positioning statement. The recorder cited by the first draft is not
part of this deposit, and no claim here depends on it.

\textbf{C.2 Projection-context congruence} (OE010-07). Following the
coalgebraic partition refinement of Birkmann et al.~(arXiv:2204.06248),
signatures \(\mathrm{sig}_\pi = F\pi \circ c\) split states until a
representation \(p\) and every context continuation agree. When
\(p(x) = p(y)\) but some context \(c\) separates them, the pair and its
separating word witness a \emph{returning projection}. This is
positioning. A related implementation exists elsewhere in the
repository, outside this deposit, and is not relied on here.

\textbf{C.3 Newton--Leibniz as a tape} (OE010-08). Integrate
\(f(x) = x\) on \([0, 10]\) with ten left-endpoint steps, exactly over
\(\mathbb{Q}\). The carrier is a path with vertices \(0, \dots, 10\).
The measured increments are \(k\) (left sum \(45\)). The exact
increments are \(d_0 F\) for \(F(k) = k^2/2\), namely
\(k + \tfrac{1}{2}\). The defect is \(\tfrac{1}{2}\) on every edge,
total \(5\). A path has no cycles, so every defect on it is exact. The
tape \texttt{C·R·E·I·T} declares the grid, registers the defect, repairs
it to \(d_0 F\), commits and seals. The integral is the telescoping sum
\(F(10) - F(0) = 50\): on a path, the discrete fundamental theorem of
calculus says that summing a coboundary returns the difference of the
potential at the endpoints. \texttt{demos.\allowbreak newton\_leibniz} computes
these numbers and the tape's receipt head;
\texttt{test\_fsm\_and\_demos.\allowbreak py::test\_newton\_leibniz\_is\_exact} pins
both. No hash is copied into this text.

\subsection{D. Spectral sequences, skills and
operads}\label{d.-spectral-sequences-skills-and-operads}

\begin{quote}
\textbf{Fence.} These are programme definitions, not theorems. No code
in this deposit implements them, and the first draft's references to
such code were dead.
\end{quote}

\begin{itemize}
\tightlist
\item
  \textbf{Self-maintenance mode} (OE010-11). A tape that re-examines its
  own detector through \texttt{C·R·Λ·I·T} would relate to the
  \(E_\infty\) page of a filtration spectral sequence. Open.
\item
  \textbf{Skill objects} (OE010-12). A \emph{skill} is proposed to be a
  residual class that survives every licensed attachment in a growth
  filtration, bound to a receipt. Definition proposal; open.
\item
  \textbf{Meta-skills} (OE010-13). A meta-skill is proposed to be a
  class that persists across a family of carriers. Open.
\item
  \textbf{Spectral sequence on a minimal complex} (OE010-14). Attaching
  a 2-cell \(\sigma\) with nonzero pairing
  \(\omega(\partial\sigma) \ne 0\) changes the question: \(\omega\) is
  no longer closed, so the class leaves \(H^1\) as an open face defect
  (\(\Gamma_2 > 0\)), and it is not repaired. Demo 3's falsifier shows
  this directly, and Theorem 5.5 shows that no lift makes a survivor
  exact under its commuting-restriction hypotheses. Support containment
  alone is insufficient: two triangle circulations +1 and −1 cancel on a
  face traversing both cycles
  (\texttt{test\_review\_regressions.\allowbreak py::test\_attaching\_a\_face\_must\_pair\_nontrivially\_with\_the\_residual}).
  A correct page-by-page account over a declared filtration is open.
\item
  \textbf{Coloured operads of skills} (OE010-15). Composition
  \(\gamma \in \mathcal{T}(c_1, \dots, c_n; d)\) with loss-ledger
  accumulation
  \(L_{\gamma \circ (\delta_i)} = L_\gamma \cup \bigcup_i \gamma^{(i)}_*(L_{\delta_i})\)
  is a design sketch. The associativity and unit laws of such an operad
  are open, and composition mints no authority.
\end{itemize}

\newpage

\section{References and their role}\label{references-and-their-role}

External literature supplies definitions and context; none of these
references authenticates a DDC runtime or supplies a licence. Local
companion links are pinned to repository commit
e4a1214bb3e47570b0f74ecc9117c98f27a26b5e.

\begin{enumerate}
\def\labelenumi{\arabic{enumi}.}
\tightlist
\item
  J. H. Carroll. \emph{Local Green, Global Red I: Evaluated Path Kernels
  for Composed Agent Reliability}, v1.2.3 (2026).
  \href{https://github.com/dbccpage/the_symbol_game/blob/e4a1214bb3e47570b0f74ecc9117c98f27a26b5e/omega_engine_atlas/coda_atlas/F004_papers/001_finite_obstruction_calculus/012_local_green_global_red_foundational_series/lggr_i_evaluated_path_kernels/lggr_i_evaluated_path_kernels.md}{Pinned
  companion source}. Context for the local/global distinction; no DDC
  proof depends on it.
\item
  J. H. Carroll. \emph{Finite Obstruction Calculus: Exact Quotient
  Diagnostics, Dual Witnesses, Two-Stage Coherence, and Transport},
  v3.2.0.
  \href{https://github.com/dbccpage/the_symbol_game/blob/e4a1214bb3e47570b0f74ecc9117c98f27a26b5e/omega_engine_atlas/coda_atlas/F005_coda/foc/coda_foc.md}{Pinned
  companion source}. Context for §3; the DDC graph formulas and
  arguments are supplied locally.
\item
  J. H. Carroll. \emph{Six jobs, not one machine: obstruction, standing,
  transport, comparison, shadow, leftover}, v3.2.0 (2026).
  \href{https://github.com/dbccpage/the_symbol_game/blob/e4a1214bb3e47570b0f74ecc9117c98f27a26b5e/omega_engine_atlas/coda_atlas/F005_coda/coda_gtla/coda_gtla_monograph.md}{Pinned
  companion source}. Background to the authorization distinction. DDC's
  refusal terminal is an explicit local policy, not an application of an
  unstated external theorem.
\item
  F. Birkmann, H.-P. Deifel and S. Milius. \emph{Distributed Coalgebraic
  Partition Refinement}. TACAS 2022, LNCS 13244, 159--177.
  \href{https://doi.org/10.1007/978-3-030-99527-0_9}{DOI
  10.1007/978-3-030-99527-0\_9},
  \href{https://arxiv.org/abs/2204.06248}{arXiv:2204.06248}. Positioning
  for Appendix C.2; no implementation is imported.
\item
  J. Meseguer. \emph{Conditional rewriting logic as a unified model of
  concurrency}. Theoretical Computer Science 96(1), 73--155 (1992).
  \href{https://doi.org/10.1016/0304-3975(92)90182-F}{DOI
  10.1016/0304-3975(92)90182-F}. Prior-art context for §16; separation
  from rewriting logic remains open.
  \href{https://maude.cs.illinois.edu/overview}{Maude project overview}
  and \href{https://maude.cs.illinois.edu/documentation}{documentation}.
  The framework can represent other operational specifications; DDC's
  chosen guards and receipt fields do not establish an expressivity
  separation.
\item
  D. Deutsch and C. Marletto. \emph{Constructor theory of information}.
  Proceedings of the Royal Society A 471, 20140540 (2015).
  \href{https://doi.org/10.1098/rspa.2014.0540}{DOI
  10.1098/rspa.2014.0540},
  \href{https://arxiv.org/abs/1405.5563}{arXiv:1405.5563}.
  Possible/impossible transformations are a comparison in §16, not a
  physics claim.
\item
  J. J. M. M. Rutten. \emph{Universal coalgebra: a theory of systems}.
  Theoretical Computer Science 249(1), 3--80 (2000).
  \href{https://ir.cwi.nl/pub/48}{CWI publication record},
  \href{https://doi.org/10.1016/S0304-3975(00)00056-6}{DOI
  10.1016/S0304-3975(00)00056-6}. Coalgebraic context for signatures and
  transition systems in Appendix C.2.
\item
  B. Coecke and A. Kissinger. \emph{Picturing Quantum Processes: A First
  Course in Quantum Theory and Diagrammatic Reasoning}. Cambridge
  University Press (2017), ISBN 978-1-107-10422-8.
  \href{https://www.cambridge.org/core/books/picturing-quantum-processes/1119568B3101F3A685BE832FEEC53E52}{Publisher
  record}. Diagrammatic context only; Appendix A's proposed adapter
  remains unproved.
\item
  L.-H. Lim. \emph{Hodge Laplacians on Graphs}. SIAM Review 62(3),
  685--715 (2020). \href{https://doi.org/10.1137/18M1223101}{DOI
  10.1137/18M1223101},
  \href{https://arxiv.org/abs/1507.05379}{arXiv:1507.05379}.
  Linear-algebra background for the decomposition stated and proved in
  §3.
\item
  The Lean project.
  \href{https://lean-lang.org/doc/reference/latest/}{Lean 4 reference
  manual}. The executable toolchain is pinned to v4.34.0 in
  lean4/lean-toolchain; the axiom report records the actual compiler
  version and commit.
\end{enumerate}

\end{document}
